%% toc: P3 39 39.7
\section{Proofs for Datalog, existential pebble games and fixed-point logics}
\label{sec:annex-p03-w03-p03-060}

Throughout, the notation is that of \cref{sec:p03-w03-p03-060-datalog-infinitary}.
Two facts are used repeatedly. \emph{Coincidence:} the truth of a formula
of \(\mathrm{L}_{\infty\omega}\) under an assignment depends only on the
values of its free variables (\cref{thm:p03-w03-p03-056-coincidence}); the
same holds for LFP, IFP and PFP formulas, by induction on formulas, because
the operator of a fixed-point formula is computed from the values of its
parameters alone. \emph{Isomorphism invariance:} an isomorphism
\(f\colon\mathbf A\to\mathbf A'\) satisfies
\(\mathbf A\models\varphi[\alpha]\iff\mathbf A'\models\varphi[f\circ\alpha]\)
for every \(\mathrm{L}_{\infty\omega}\)-formula, by induction on formulas
(the atomic case is the definition of isomorphism, and \(f\) is a
bijection in the quantifier case).

\begin{annexitem}{annex:p03-w03-p03-060-existential-game}{thm:p03-w03-p03-060-existential-game}{Existential pebble games characterize existential \(k\)-variable preservation}
\begin{proof}
We treat the \((\exists,k)\) version and state at each step what changes in
the \((\exists,\neq,k)\) version.

\emph{Step 0: good positions and atomic formulas.} Let \((\alpha,\beta)\) be
a position with domain \(D\). For a term \(t\) whose variables have labels
in \(D\), the pair \((t^{\mathbf A}[\alpha],t^{\mathbf B}[\beta])\) lies in
\(\rho_{\alpha,\beta}\): for \(t=v_i\) it is \((\alpha(i),\beta(i))\), and
for a constant \(c\) it is \((c^{\mathbf A},c^{\mathbf B})\).

(a) Suppose \((\alpha,\beta)\) is \(\exists\)-good and \(\varphi\) is atomic
with \(\operatorname{fr}(\varphi)\subseteq D\) and
\(\mathbf A\models\varphi[\alpha]\). If \(\varphi\) is \(t=t'\), the two
values in \(\mathbf A\) coincide; since \(\rho=\rho_{\alpha,\beta}\) is a
function, their images \(t^{\mathbf B}[\beta]\) and \(t'^{\mathbf B}[\beta]\)
coincide, so \(\mathbf B\models\varphi[\beta]\). If \(\varphi\) is
\(R(t_1,\ldots,t_r)\), then \(R^{\mathbf A}\) holds of the values
\(t_j^{\mathbf A}[\alpha]\in\dom\rho\), and since \(\rho\) is a partial
homomorphism, \(R^{\mathbf B}\) holds of their images
\(t_j^{\mathbf B}[\beta]\).

(b) In the \((\neq)\) version, if moreover \(\rho\) is one-to-one and
\(\mathbf A\models t\neq t'[\alpha]\), the two values in \(\mathbf A\) are
distinct, hence so are their images under \(\rho\), and
\(\mathbf B\models t\neq t'[\beta]\).

(c) Conversely, suppose \((\alpha,\beta)\) is not \(\exists\)-good. Every
element of \(\dom\rho\) is the value in \(\mathbf A\) of a term with labels
in \(D\). If \(\rho\) is not a function, there are terms \(t,t'\) with equal
values in \(\mathbf A\) and different values in \(\mathbf B\); then
\(t=t'\) is true at \(\alpha\) and false at \(\beta\). If \(\rho\) is a
function but not a partial homomorphism (it sends constants to constants by
construction), some \(R^{\mathbf A}(a_1,\ldots,a_r)\) with
\(a_j\in\dom\rho\) fails for the images; writing \(a_j\) as the value of a
term \(t_j\), the atom \(R(t_1,\ldots,t_r)\) is true at \(\alpha\) and false
at \(\beta\). In the \((\neq)\) version, if \(\rho\) is a partial
homomorphism but not one-to-one, two terms have different values in
\(\mathbf A\) and equal values in \(\mathbf B\), and \(t\neq t'\) is true at
\(\alpha\) and false at \(\beta\). In every case the formula found is an
\(\exists\mathrm{FO}^k\)-formula (an \(\exists\mathrm{FO}^k(\neq)\)-formula)
with free labels in \(D\).

\emph{(1)\(\Rightarrow\)(2).} Let \(\mathcal W\) be the set of positions
from which the Duplicator has a winning strategy; it contains
\((\alpha,\beta)\). Each member is \(\exists\)-good, because the initial
position of a play counts. If \(\Sigma\) wins from \(p\in\mathcal W\) and the
Spoiler makes a removal or a placement \((i,a)\) to which \(\Sigma\) answers
\(b\), the tail strategy wins from the resulting position, as recorded in
\cref{def:p03-w03-p03-056-pebble-game}. Hence \(\mathcal W\) is closed
under removal and has the extension property, and is an \(\exists\)-family.

\emph{(2)\(\Rightarrow\)(1).} Let \(\mathcal W\) be an \(\exists\)-family
containing \((\alpha,\beta)\). The Duplicator plays as follows: at a
position in \(\mathcal W\), a removal leads to a position in
\(\mathcal W\); to a placement \((i,a)\) she answers with some \(b\) such
that the new position lies in \(\mathcal W\), which exists by the extension
property. By induction on the number of rounds every position of the play
lies in \(\mathcal W\) and is therefore \(\exists\)-good, so the strategy
wins.

\emph{(2)\(\Rightarrow\)(3).} Let \(\mathcal W\) be an \(\exists\)-family.
We show by induction on the formation of
\(\varphi\in\exists\mathrm{L}^k_{\infty\omega}\): for every
\((\alpha',\beta')\in\mathcal W\) with
\(\operatorname{fr}(\varphi)\subseteq\dom\alpha'\),
\(\mathbf A\models\varphi[\alpha']\) implies
\(\mathbf B\models\varphi[\beta']\). Atomic formulas: Step~0(a). For
\(\bigwedge\Phi\) and \(\bigvee\Phi\), each member of \(\Phi\) has free
labels in \(\dom\alpha'\), and the induction hypothesis applied to all
members (for \(\bigwedge\)) or to one true member (for \(\bigvee\)) gives
the claim. For \(\exists v_i\psi\): choose \(a\in A\) with
\(\mathbf A\models\psi[\alpha'[i\mapsto a]]\); the extension property gives
\(b\) with \((\alpha'[i\mapsto a],\beta'[i\mapsto b])\in\mathcal W\); since
\(\operatorname{fr}(\psi)\subseteq\operatorname{fr}(\exists v_i\psi)\cup\{i\}\),
the induction hypothesis gives \(\mathbf B\models\psi[\beta'[i\mapsto b]]\),
so \(\mathbf B\models\exists v_i\psi[\beta']\). Applied to
\((\alpha,\beta)\) this is (3). In the \((\neq)\) version the additional
basic formulas \(t\neq t'\) are handled by Step~0(b).

\emph{(3)\(\Rightarrow\)(2).} Let \(\mathcal W\) be the set of positions
\((\alpha',\beta')\) over \([k]\) with
\((\mathbf A,\alpha')\preceq^{\exists,k}_{\infty\omega}(\mathbf B,\beta')\);
it contains \((\alpha,\beta)\). Every member is \(\exists\)-good: otherwise
Step~0(c) produces a formula of \(\exists\mathrm{L}^k_{\infty\omega}\) with
free labels in the domain, true at \(\alpha'\) and false at \(\beta'\).
Removal shrinks the set of formulas to be preserved, so \(\mathcal W\) is
closed under removal. For the extension property, let
\((\alpha',\beta')\in\mathcal W\) with domain \(D'\), \(i\in[k]\) and
\(a\in A\), and suppose that no \(b\in B\) satisfies
\((\alpha'[i\mapsto a],\beta'[i\mapsto b])\in\mathcal W\). Then for every
\(b\in B\) we may choose \(\varphi_b\in\exists\mathrm{L}^k_{\infty\omega}\)
with \(\operatorname{fr}(\varphi_b)\subseteq D'\cup\{i\}\),
\(\mathbf A\models\varphi_b[\alpha'[i\mapsto a]]\) and
\(\mathbf B\not\models\varphi_b[\beta'[i\mapsto b]]\). Put
\(\varphi\defeq\exists v_i\bigwedge\{\varphi_b:b\in B\}\), a formula of
\(\exists\mathrm{L}^k_{\infty\omega}\) with free labels in \(D'\). The
witness \(a\) shows \(\mathbf A\models\varphi[\alpha']\), so
\(\mathbf B\models\varphi[\beta']\): some \(b_0\in B\) satisfies every
\(\varphi_b\) at \(\beta'[i\mapsto b_0]\), in particular \(\varphi_{b_0}\),
contrary to its choice. In the \((\neq)\) version all formulas are taken
from \(\exists\mathrm{L}^k_{\infty\omega}(\neq)\), Step~0(c) supplies the
inequality, and \(\varphi\) again belongs to that language.

\emph{First-order part.} Since \(\exists\mathrm{FO}^k\subseteq\exists\mathrm{L}^k_{\infty\omega}\),
(3) implies
\((\mathbf A,\alpha)\preceq^{\exists,k}_{\omega\omega}(\mathbf B,\beta)\).
If \(\mathbf B\) is finite and this relation holds, repeat the proof of
(3)\(\Rightarrow\)(2) with \(\exists\mathrm{FO}^k\) in place of
\(\exists\mathrm{L}^k_{\infty\omega}\): Step~0(c) yields first-order
formulas, and the conjunction \(\bigwedge\{\varphi_b:b\in B\}\) is finite,
so \(\varphi\in\exists\mathrm{FO}^k\). This gives (2) and hence (1)--(3).
The \((\neq)\) version is the same with \(\exists\mathrm{FO}^k(\neq)\).
\end{proof}
\end{annexitem}

\begin{annexitem}{annex:p03-w03-p03-060-existential-definability}{thm:p03-w03-p03-060-existential-definability}{Existential definability on classes of finite structures}
\begin{proof}
Suppose \(\varphi\in\exists\mathrm{L}^k_{\infty\omega}\) is a sentence
defining \(\mathcal K\) on \(\mathcal C\). If \(\mathbf A\in\mathcal K\),
\(\mathbf B\in\mathcal C\) and the Duplicator wins the
\((\exists,k)\)-pebble game on them from the empty position, then
\(\mathbf A\preceq^{\exists,k}_{\infty\omega}\mathbf B\)
(\cref{thm:p03-w03-p03-060-existential-game}, with \(D=\emptyset\)); as
\(\mathbf A\models\varphi\), also \(\mathbf B\models\varphi\), so
\(\mathbf B\in\mathcal K\).

Conversely, fix \(k\) with the stated property. Every finite
\(\sigma\)-structure is isomorphic to one whose universe is
\(\{1,\ldots,n\}\) for some \(n\ge1\), and these form a set; let
\(\mathcal R\) be the set of such structures that lie in \(\mathcal C\) up
to isomorphism, taken with one representative per isomorphism type and
each representative in \(\mathcal C\). For
\(\mathbf A\in\mathcal R\cap\mathcal K\) and
\(\mathbf B\in\mathcal R\setminus\mathcal K\), the Duplicator does not win
the \((\exists,k)\)-pebble game on \(\mathbf A\) and \(\mathbf B\) (else
\(\mathbf B\in\mathcal K\)), so by
\cref{thm:p03-w03-p03-060-existential-game} there is a sentence
\(\varphi_{\mathbf A,\mathbf B}\in\exists\mathrm{L}^k_{\infty\omega}\) true
in \(\mathbf A\) and false in \(\mathbf B\). Put
\[
 \Phi\defeq\bigvee_{\mathbf A\in\mathcal R\cap\mathcal K}\ \bigwedge_{\mathbf B\in\mathcal R\setminus\mathcal K}\varphi_{\mathbf A,\mathbf B},
\]
a sentence of \(\exists\mathrm{L}^k_{\infty\omega}\). Let
\(\mathbf C\in\mathcal C\). If \(\mathbf C\in\mathcal K\), it is isomorphic
to some \(\mathbf A\in\mathcal R\), which lies in \(\mathcal K\) because
\(\mathcal K\) is closed under isomorphism; \(\mathbf A\) satisfies the
disjunct indexed by \(\mathbf A\), so \(\mathbf C\models\Phi\) by
isomorphism invariance. If \(\mathbf C\notin\mathcal K\), it is isomorphic
to some \(\mathbf B\in\mathcal R\setminus\mathcal K\), and every disjunct
contains the conjunct \(\varphi_{\mathbf A,\mathbf B}\), false in
\(\mathbf B\); so \(\mathbf B\not\models\Phi\) and
\(\mathbf C\not\models\Phi\). Thus \(\Phi\) defines \(\mathcal K\) on
\(\mathcal C\). Both directions use
\cref{thm:p03-w03-p03-060-existential-game} only through the equivalence
of items~1 and~3, which holds for the \((\exists,\neq,k)\) game and
\(\exists\mathrm{L}^k_{\infty\omega}(\neq)\); the same two arguments with
these replacements prove the \((\neq)\) version.
\end{proof}
\end{annexitem}

\begin{annexitem}{annex:p03-w03-p03-060-preservation}{thm:p03-w03-p03-060-preservation}{Preservation under homomorphisms}
\begin{proof}
Choose \(k\) with \(\varphi\in\exists\mathrm{L}^k_{\infty\omega}\). Let
\(\mathcal W\) be the set of all positions \((\alpha',h\circ\alpha')\) with
\(\dom\alpha'\subseteq[k]\). For such a position,
\(\rho_{\alpha',h\circ\alpha'}\) is contained in the graph of \(h\),
because \(h(c^{\mathbf A})=c^{\mathbf B}\); hence it is a function, and,
as a restriction of a homomorphism, a partial homomorphism, one-to-one if
\(h\) is. Removal preserves the form of the position, and a placement
\((i,a)\) is answered by \(b=h(a)\). So \(\mathcal W\) is an
\(\exists\)-family, and an \(\exists{\neq}\)-family if \(h\) is one-to-one.
It contains \((\alpha,h\circ\alpha)\), and
\cref{thm:p03-w03-p03-060-existential-game} ((2)\(\Rightarrow\)(3)) gives
the claim.
\end{proof}
\end{annexitem}

\begin{annexitem}{annex:p03-w03-p03-060-cliques}{thm:p03-w03-p03-060-cliques}{The existential game is not symmetric}
\begin{proof}
\emph{Spoiler on \(\mathbf K_{k+1}\) and \(\mathbf K_k\).} Let the Spoiler
place pebbles \(1,\ldots,k+1\), one per round, on the \(k+1\) distinct
vertices \(a_1,\ldots,a_{k+1}\) of \(\mathbf K_{k+1}\). Whatever the
Duplicator answers, after round \(k+1\) the answers
\(b_1,\ldots,b_{k+1}\) lie in a set of \(k\) vertices, so \(b_i=b_j\) for
some \(i\neq j\). Then \(E(a_i,a_j)\) holds in \(\mathbf K_{k+1}\) while
\(E(b_i,b_j)=E(b_i,b_i)\) fails in \(\mathbf K_k\), which has no loops;
the final position is not \(\exists\)-good. So no strategy of the
Duplicator is winning.

\emph{Duplicator on \(\mathbf K_k\) and \(\mathbf K_{k+1}\).} Fix an
injective map \(e\) from the vertices of \(\mathbf K_k\) to those of
\(\mathbf K_{k+1}\). If \(E(a,a')\) holds in \(\mathbf K_k\), then
\(a\neq a'\), so \(e(a)\neq e(a')\) and \(E(e(a),e(a'))\) holds in
\(\mathbf K_{k+1}\): \(e\) is a one-to-one homomorphism. For every
\(j\ge1\), the positions \((\alpha,e\circ\alpha)\) with
\(\dom\alpha\subseteq[j]\) form an \(\exists{\neq}\)-family by the argument
in the proof of \cref{thm:p03-w03-p03-060-preservation}, and it contains
the empty position. By \cref{thm:p03-w03-p03-060-existential-game}
((2)\(\Rightarrow\)(1)) the Duplicator wins the \((\exists,\neq,j)\)-pebble
game, and a fortiori the \((\exists,j)\)-pebble game, whose goodness
condition is weaker.
\end{proof}
\end{annexitem}

\begin{annexitem}{annex:p03-w03-p03-060-substitution}{thm:p03-w03-p03-060-substitution}{Substituting a defined relation for a relation symbol}
\begin{proof}
Write \(\theta_l^{\mathbf A,s}\) as in the statement and let
\(\mathbf A^\ast\defeq(\mathbf A,\theta_1^{\mathbf A,s},\ldots,\theta_n^{\mathbf A,s})\).
If \(s'\) agrees with \(s\) on the parameters, then
\(\theta_l^{\mathbf A,s'}=\theta_l^{\mathbf A,s}\) by coincidence, because
the free variables of \(\theta_l\) are among \(x_1,\ldots,x_{r_l}\) (which
are overwritten) and the parameters.

We prove by induction on the formation of subformulas \(\chi'\) of \(\chi\)
that \(\mathbf A\models\chi'[\bar R:=\bar\theta][s']\iff\mathbf A^\ast\models\chi'[s']\)
for every assignment \(s'\) to the free variables of \(\chi'\) and the
parameters that agrees with \(s\) on the parameters.

\emph{Atom \(R_l(t_1,\ldots,t_{r_l})\).} No \(z_j\) occurs in \(\chi\), so
the terms \(t_j\) are not among \(z_1,\ldots,z_{r_l}\), and the outer quantifier group
\(\exists\bar z(\bigwedge_jz_j=t_j\wedge\cdots)\) holds exactly when the
rest holds with \(z_j\) set to \(d_j\defeq t_j^{\mathbf A}[s']\). The second
group then holds exactly when \(\theta_l\) holds at
\(s'[\bar z\mapsto\bar d][\bar x\mapsto\bar d]\). No \(z_j\) is free in
\(\theta_l\), so by coincidence this is truth of \(\theta_l\) at
\(s'[\bar x\mapsto\bar d]\), that is,
\(\bar d\in\theta_l^{\mathbf A,s'}=\theta_l^{\mathbf A,s}\), which is
\(\mathbf A^\ast\models R_l(t_1,\ldots,t_{r_l})[s']\).

\emph{Other atoms.} They are unchanged, and \(\mathbf A\), \(\mathbf A^\ast\)
agree on \(\tau\).

\emph{\(\neg\), \(\bigwedge\), \(\bigvee\).} Substitution commutes with these
constructors, and the induction hypothesis applies to the immediate
subformulas at the same \(s'\).

\emph{\(\exists y\,\chi''\) and \(\forall y\,\chi''\).} The variable \(y\)
is bound in \(\chi\), so it is not a parameter; hence \(s'[y\mapsto a]\)
agrees with \(s\) on the parameters for every \(a\), and the induction
hypothesis applies to \(\chi''\) at each \(s'[y\mapsto a]\).

Taking \(\chi'=\chi\) and \(s'=s\) proves the equivalence. The variable
claim is read off the replacement. The replacement uses only \(\exists\),
\(\bigwedge\) and equality atoms, so it stays in
\(\exists\mathrm{L}_{\infty\omega}\) (with \(\neq\)) when \(\chi\) and the
\(\theta_l\) do. Finally, each copy of \(\theta_l\) is placed, under
\(\exists\) and \(\bigwedge\) only, at the position of an occurrence of
\(R_l\); if that occurrence is positive and \(P\) is positive in
\(\theta_l\), the occurrences of \(P\) in the copy are positive, and the
occurrences of \(P\) in \(\chi\) keep their polarity.
\end{proof}
\end{annexitem}

\begin{annexitem}{annex:p03-w03-p03-060-fp-in-lomega}{thm:p03-w03-p03-060-fp-in-lomega}{LFP, IFP and PFP are contained in \(\mathrm{L}^\omega_{\infty\omega}\) on finite structures}
\begin{proof}
\emph{Renaming.} Let \(\xi\) be an \(\mathrm{L}_{\infty\omega}\)-formula
with finitely many variables, \(y\) a variable and \(y'\) a variable not
occurring in \(\xi\); let \(\xi^{y\to y'}\) replace every occurrence of
\(y\), free or bound, by \(y'\). By induction on \(\xi\),
\(\mathbf A\models\xi^{y\to y'}[s]\iff\mathbf A\models\xi[\tilde s]\),
where \(\tilde s\) sends \(y\) to \(s(y')\) and agrees with \(s\)
elsewhere: atoms are evaluated at the same values; the connectives commute;
for \(\exists y\) the quantifier becomes \(\exists y'\) and
\(\widetilde{s[y'\mapsto a]}\) agrees with \(\tilde s[y\mapsto a]\) on the
variables of \(\xi\); for \(\exists x\) with \(x\neq y\) we have
\(x\neq y'\) and the two updates commute; \(\forall\) is the same. Hence
replacing a subformula \(Qy\,\xi\) by \(Qy'\,\xi^{y\to y'}\), with \(y'\)
new, does not change the truth value of the whole formula under any
assignment and does not change its free variables.

\emph{Construction.} By induction on \(\varphi\) we build
\(\varphi^\ast\in\mathrm{L}_{\infty\omega}\) with finitely many variables
and free variables among those of \(\varphi\), equivalent to \(\varphi\) on
every finite structure interpreting its free relation variables. Atomic
formulas are their own translations; \(\neg\), \(\wedge\), \(\vee\),
\(\exists\), \(\forall\) are applied to the translations of the finitely
many immediate subformulas.

Let \(\varphi=[\mathbf{op}\,R\bar x\,.\,\psi](\bar t)\) with
\(\mathbf{op}\in\{\mathbf{lfp},\mathbf{gfp},\mathbf{ifp},\mathbf{pfp}\}\),
\(\bar x=(x_1,\ldots,x_r)\). By the induction hypothesis \(\psi^\ast\)
exists. Let \(\bar p\) be the free variables of \(\psi^\ast\) outside
\(\bar x\); they are among the parameters of \(\varphi\). For each
\(p\in\bar p\) choose a new variable \(p'\) and, using Renaming, replace
each maximal subformula \(Qp\,\xi\) of \(\psi^\ast\) by
\(Qp'\,\xi^{p\to p'}\). The result \(\psi^\circ\) is equivalent to
\(\psi^\ast\), has the same free variables, and binds no variable of
\(\bar p\). Choose new variables \(z_1,\ldots,z_r\) and define
\[
 \psi^0\defeq\begin{cases}\top&\mathbf{op}=\mathbf{gfp},\\ \bot&\text{otherwise,}\end{cases}
 \qquad
 \psi^{m+1}\defeq\begin{cases}
 (R(\bar x)\vee\psi^\circ)[R:=\psi^m]&\mathbf{op}=\mathbf{ifp},\\
 \psi^\circ[R:=\psi^m]&\text{otherwise,}\end{cases}
\]
with the substitution of \cref{thm:p03-w03-p03-060-substitution} for
\(w=r\). By induction on \(m\): the free variables of \(\psi^m\) are among
\(\bar x\) and \(\bar p\) (the replacement of an atom binds \(\bar x\)
around the copy of \(\psi^m\)), so the parameters of \(\psi^m\) are among
\(\bar p\) and are bound nowhere in \(\psi^\circ\), no \(z_j\) is free in
\(\psi^m\), and no \(z_j\) occurs in \(\psi^\circ\); thus the lemma
applies at every \(m\), and the variables of \(\psi^m\) lie in the finite
set \(V\) of variables of \(\psi^\circ\) together with \(\bar x\) and
\(\bar z\), independently of \(m\).

Fix a finite structure \(\mathbf A\) interpreting the free relation
variables, an assignment \(s\) to the free variables of \(\varphi\)
(extended arbitrarily to \(\bar p\) if needed; the parameters of
\(\varphi\) include \(\bar p\)), and let
\(F(S)=\{\bar a:(\mathbf A,S)\models\psi[s[\bar x\mapsto\bar a]]\}\), which
by the induction hypothesis and Renaming is also defined by \(\psi^\circ\).
Let \(R^m\defeq\{\bar a:\mathbf A\models\psi^m[s[\bar x\mapsto\bar a]]\}\).
Then \(R^0=\varnothing\) (\(A^r\) for \(\mathbf{gfp}\)), and by the lemma,
applied at the assignment \(s[\bar x\mapsto\bar a]\), which agrees with
\(s\) on \(\bar p\), \(R^{m+1}=F(R^m)\), and
\(R^{m+1}=R^m\cup F(R^m)\) for \(\mathbf{ifp}\). Again by the lemma, with
\(\chi=R(\bar t)\) (whose variables are those of \(\bar t\), none a
\(z_j\)), the formula \(\eta_m\defeq R(\bar t)[R:=\psi^m]\) holds at \(s\)
exactly when \(\bar t^{\mathbf A}[s]\in R^m\).

\(\mathbf{lfp}\): \(\varphi\) is an LFP formula, so \(R\) is positive in
\(\psi\) and \(F\) is monotone
(\cref{thm:p03-w03-p03-059-positive-monotone}). Then \(R^0\subseteq R^1\),
and \(R^m\subseteq R^{m+1}\) gives
\(R^{m+1}=F(R^m)\subseteq F(R^{m+1})=R^{m+2}\). An increasing sequence of
subsets of the finite set \(A^r\) is constant from some \(m_0\) on, so
\(R^{m_0}\) is a fixed point. If \(F(S)=S\), then \(R^0\subseteq S\) and
\(R^m\subseteq S\) gives \(R^{m+1}=F(R^m)\subseteq F(S)=S\); so
\(R^{m_0}=\bigcup_mR^m\) is the least fixed point. Put
\(\varphi^\ast\defeq\bigvee_{m\in\N}\eta_m\).

\(\mathbf{gfp}\): with \(R^0=A^r\), \(R^1\subseteq R^0\) and monotonicity
give a decreasing sequence, constant from some \(m_0\) on; \(R^{m_0}\) is a
fixed point, and every fixed point \(S\) satisfies \(S\subseteq R^m\) for
all \(m\) by the dual induction, so \(R^{m_0}=\bigcap_mR^m\) is the
greatest fixed point. Put \(\varphi^\ast\defeq\bigwedge_{m\in\N}\eta_m\).

\(\mathbf{ifp}\): the inflationary fixed point is \(\bigcup_mR^m\) by
\cref{def:p03-w03-p03-060-ifp}. Put \(\varphi^\ast\defeq\bigvee_{m\in\N}\eta_m\).

\(\mathbf{pfp}\): put
\[
 \varphi^\ast\defeq\bigvee_{m\in\N}\Bigl(\forall x_1\cdots\forall x_r
 \bigl((\neg\psi^m\vee\psi^{m+1})\wedge(\neg\psi^{m+1}\vee\psi^m)\bigr)\wedge\eta_m\Bigr).
\]
The first conjunct of the \(m\)-th disjunct holds at \(s\) exactly when
\(R^m=R^{m+1}\). If some stage is a fixed point, the partial fixed point is
\(R^m\) for every \(m\) with \(R^m=R^{m+1}\)
(\cref{def:p03-w03-p03-059-pfp}), so \(\varphi^\ast\) holds exactly when
\(\bar t^{\mathbf A}[s]\) lies in the partial fixed point. If no stage is a
fixed point, every disjunct is false and the partial fixed point is empty,
so \(\varphi\) is false as well.

In each case the variables of \(\varphi^\ast\) lie in
\(V\cup\operatorname{var}(\bar t)\), a finite set, and its free variables
are among those of \(\bar t\) and \(\bar p\), hence among those of
\(\varphi\). This completes the induction. Counting the construction gives
an explicit bound: if \(\psi^\ast\) has \(k_\psi\) variables, \(\varphi^\ast\)
has at most \(k_\psi+|\bar p|+2r+|\operatorname{var}(\bar t)|\) variables.
\end{proof}
\end{annexitem}

\begin{annexitem}{annex:p03-w03-p03-060-variable-substitution}{thm:p03-w03-p03-060-variable-substitution}{Substituting terms for variables in \(k\)-variable formulas}
\begin{proof}
For a map \(\rho\colon[k]\to T_k\) and a term \(t\) put
\(t\rho\defeq\rho(i)\) if \(t=v_i\) and \(t\rho\defeq t\) if \(t\) is a
constant symbol. Then \((t\rho)^{\mathbf A}[\alpha]=t^{\mathbf A}[\rho^\alpha]\):
for \(t=v_i\) both sides are \(\rho(i)^{\mathbf A}[\alpha]\), and for a
constant both are \(c^{\mathbf A}\). For a set \(I\subseteq[k]\) write
\(V(\rho,I)\defeq\{j:v_j\text{ occurs in }\rho(i)\text{ for some }i\in I\}\);
since a term contains at most one variable, \(|V(\rho,I)|\le|I|\).

We define \(\varphi_\rho\) by recursion on the formation of \(\varphi\),
for all \(\rho\) at once, and prove the equivalence for all \(\rho\) and
all \(\alpha\colon[k]\to A\), together with
\(\operatorname{fr}(\varphi_\rho)\subseteq V(\rho,\operatorname{fr}(\varphi))\).

\emph{Atomic formulas and inequalities.}
\(R(t_1,\ldots,t_r)_\rho\defeq R(t_1\rho,\ldots,t_r\rho)\),
\((t=t')_\rho\defeq(t\rho=t'\rho)\) and
\((t\neq t')_\rho\defeq(t\rho\neq t'\rho)\). Each holds at \(\alpha\)
exactly when the original holds at \(\rho^\alpha\), by the identity for
term values. The variables of \(t\rho\) are those of \(\rho(i)\) for
\(v_i\) occurring in \(t\), so the free labels lie in
\(V(\rho,\operatorname{fr}(\varphi))\).

\emph{\(\neg\), \(\bigwedge\), \(\bigvee\).}
\((\neg\psi)_\rho\defeq\neg(\psi_\rho)\),
\((\bigwedge\Phi)_\rho\defeq\bigwedge\{\psi_\rho:\psi\in\Phi\}\) and
\((\bigvee\Phi)_\rho\defeq\bigvee\{\psi_\rho:\psi\in\Phi\}\). The
equivalence and the bound on free labels pass through each constructor,
and the new set is finite when \(\Phi\) is.

\emph{\(\exists v_i\psi\) and \(\forall v_i\psi\).} Let
\(U\defeq V(\rho,\operatorname{fr}(\psi)\setminus\{i\})\). Then
\(|U|\le k-1\), so \([k]\setminus U\) is not empty; let \(j\) be its least
element and \(\rho'\defeq\rho[i\mapsto v_j]\). Define
\((\exists v_i\psi)_\rho\defeq\exists v_j\,\psi_{\rho'}\) and
\((\forall v_i\psi)_\rho\defeq\forall v_j\,\psi_{\rho'}\). For \(a\in A\),
by the induction hypothesis
\(\mathbf A\models\psi_{\rho'}[\alpha[j\mapsto a]]\) iff
\(\mathbf A\models\psi[\rho'^{\alpha[j\mapsto a]}]\). The assignment
\(\rho'^{\alpha[j\mapsto a]}\) sends \(i\) to the value of \(v_j\), which
is \(a\), and sends each \(l\in\operatorname{fr}(\psi)\setminus\{i\}\) to
\(\rho(l)^{\mathbf A}[\alpha[j\mapsto a]]=\rho(l)^{\mathbf A}[\alpha]=\rho^\alpha(l)\),
because \(v_j\) does not occur in \(\rho(l)\). So it agrees with
\(\rho^\alpha[i\mapsto a]\) on \(\operatorname{fr}(\psi)\), and by
\cref{thm:p03-w03-p03-056-coincidence}
\(\mathbf A\models\psi_{\rho'}[\alpha[j\mapsto a]]\) iff
\(\mathbf A\models\psi[\rho^\alpha[i\mapsto a]]\). Taking ``some \(a\)''
gives the claim for \(\exists\), and ``every \(a\)'' for \(\forall\). For
the free labels,
\(\operatorname{fr}(\psi_{\rho'})\subseteq V(\rho',\operatorname{fr}(\psi))\subseteq\{j\}\cup U\),
so
\(\operatorname{fr}(\exists v_j\psi_{\rho'})\subseteq U=V(\rho,\operatorname{fr}(\exists v_i\psi))\).

\emph{Membership in \(\mathcal L\).} \(\varphi_\rho\) is formed with the
constructors used in forming \(\varphi\), with sets of the same or smaller
size, and its basic formulas are of the same kinds (atoms, equalities, and
inequalities only where \(\varphi\) has them). So \(\varphi_\rho\) lies in
each of the six languages that contains \(\varphi\). In the example of the
body text, \(U=V(\rho,\{1\})=\{2\}\), \(j=1\), \(\rho'(1)=v_2\),
\(\rho'(2)=v_1\), and \(E(v_1,v_2)_{\rho'}=E(v_2,v_1)\).
\end{proof}
\end{annexitem}

\begin{annexitem}{annex:p03-w03-p03-060-kv-stages}{thm:p03-w03-p03-060-kv-stages}{Kolaitis and Vardi: the stages of a \(k\)-variable formula are \(k\)-variable formulas}
\begin{proof}
\emph{Normal form.} Let \(W\) be the declared supply of \(k\ge1\)
variables, containing every variable of \(\varphi\) and all its coordinate
variables, and let \(\pi\colon W\to\{v_1,\ldots,v_k\}\) be a bijection with
\(\pi(x_j)=v_j\) (\(j\le m\); these distinct variables belong to \(W\),
so \(m\le k\)). By
\cref{thm:p03-w03-p03-056-renaming}, \(\varphi^\pi\) is an
\(\mathrm{FO}^k_{\omega\omega}\)-formula of vocabulary \(\sigma\cup\{S\}\)
with free labels in \([m]\), and it holds at \(\alpha\) exactly when
\(\varphi\) holds at the assignment \(x_j\mapsto\alpha(j)\); so it has the
same operator \(\Phi\) on every structure. We write \(\varphi\) for
\(\varphi^\pi\).

\emph{Replacing \(S\).} For \(\theta\in\mathrm{L}^k_{\infty\omega}(\sigma)\)
with free labels in \([m]\) put
\(\theta^{\mathbf A}\defeq\{(\alpha(1),\ldots,\alpha(m)):\alpha\colon[k]\to A,\ \mathbf A\models\theta[\alpha]\}\);
by \cref{thm:p03-w03-p03-056-coincidence}, \(\mathbf A\models\theta[\alpha]\)
iff \((\alpha(1),\ldots,\alpha(m))\in\theta^{\mathbf A}\). For a formula
\(\chi\) of vocabulary \(\sigma\cup\{S\}\) over \(v_1,\ldots,v_k\) let
\(\chi[S:=\theta]\) replace each atomic subformula \(S(t_1,\ldots,t_m)\) by
\(\theta_{\rho_{\bar t}}\), where \(\rho_{\bar t}(j)\defeq t_j\) for
\(j\le m\) and \(\rho_{\bar t}(j)\defeq v_j\) for \(j>m\)
(\cref{thm:p03-w03-p03-060-variable-substitution}). We claim: for every
\(\alpha\colon[k]\to A\),
\(\mathbf A\models\chi[S:=\theta][\alpha]\iff(\mathbf A,\theta^{\mathbf A})\models\chi[\alpha]\).
By induction on \(\chi\), for all \(\alpha\): at an atom \(S(\bar t)\), the
lemma gives \(\mathbf A\models\theta_{\rho_{\bar t}}[\alpha]\) iff
\(\mathbf A\models\theta[\rho_{\bar t}^\alpha]\) iff
\((t_1^{\mathbf A}[\alpha],\ldots,t_m^{\mathbf A}[\alpha])\in\theta^{\mathbf A}\),
which is the truth of \(S(\bar t)\) in \((\mathbf A,\theta^{\mathbf A})\);
other atoms are unchanged; the connectives commute with the replacement;
and for \(\exists v_i\) and \(\forall v_i\) the induction hypothesis is used
at every \(\alpha[i\mapsto a]\). The free labels of \(\theta_{\rho_{\bar t}}\)
are among those of the variables of \(t_1,\ldots,t_m\), that is among
\(\operatorname{fr}(S(\bar t))\), so \(\chi[S:=\theta]\) has no free label
outside \(\operatorname{fr}(\chi)\). If \(\chi\) is first-order and
\(\theta\in\mathrm{FO}^k_{\omega\omega}\), then
\(\chi[S:=\theta]\in\mathrm{FO}^k_{\omega\omega}\), since finitely many
atoms are replaced by first-order formulas.

\emph{Item 1.} Let \(\bot\defeq\bigvee\emptyset\) and
\(\top\defeq\bigwedge\emptyset\), closed first-order formulas false and
true at every assignment. Their sets of free labels are empty, including
when \(m=0\). Put \(\varphi^0\defeq\bot\),
\(\varphi^{n+1}\defeq\varphi[S:=\varphi^n]\);
\(\varphi^0_{\mathrm{inf}}\defeq\bot\),
\(\varphi^{n+1}_{\mathrm{inf}}\defeq(S(v_1,\ldots,v_m)\vee\varphi)[S:=\varphi^n_{\mathrm{inf}}]\);
\(\varphi_0\defeq\top\), \(\varphi_{n+1}\defeq\varphi[S:=\varphi_n]\). The
formula \(S(v_1,\ldots,v_m)\vee\varphi\) is in
\(\mathrm{FO}^k_{\omega\omega}(\sigma\cup\{S\})\) with free labels in
\([m]\). By induction on \(n\), all these formulas are in
\(\mathrm{FO}^k_{\omega\omega}(\sigma)\) with free labels in \([m]\), and by
the claim \((\varphi^{n+1})^{\mathbf A}=\Phi((\varphi^n)^{\mathbf A})\),
\((\varphi^{n+1}_{\mathrm{inf}})^{\mathbf A}=(\varphi^n_{\mathrm{inf}})^{\mathbf A}\cup\Phi((\varphi^n_{\mathrm{inf}})^{\mathbf A})\)
and \((\varphi_{n+1})^{\mathbf A}=\Phi((\varphi_n)^{\mathbf A})\), with
\((\varphi^0)^{\mathbf A}=\varnothing\) and
\((\varphi_0)^{\mathbf A}=A^m\). So the formulas define \(\Phi^n\),
\(\Phi^n_{\mathrm{inf}}\) and \(\Phi_n\) on every structure; no finiteness
was used.

\emph{Item 2.} Let \(\mathbf A\) be finite. Partial fixed point: the
\(n\)-th disjunct holds at \(\alpha\) exactly when
\((\alpha(1),\ldots,\alpha(m))\in\Phi^n\) and \(\Phi^n=\Phi^{n+1}\), since
the universal block ranges over all values of the labels \(1,\ldots,m\),
which are the only free labels of \(\varphi^n\) and \(\varphi^{n+1}\). If
\(\Phi^n=\Phi^{n+1}\) for some \(n\), the stages are constant from \(n\) on,
all such \(n\) give the same set, and that set is the partial fixed point
(\cref{def:p03-w03-p03-059-pfp}); otherwise every disjunct is false and the
partial fixed point is empty. Inflationary fixed point: it is
\(\bigcup_n\Phi^n_{\mathrm{inf}}\) by \cref{def:p03-w03-p03-060-ifp}.
Least and greatest fixed points: if \(S\) is positive in \(\varphi\), then
\(\Phi\) is monotone (\cref{thm:p03-w03-p03-059-positive-monotone}), and the
cases \(\mathbf{lfp}\) and \(\mathbf{gfp}\) of the proof of
\cref{thm:p03-w03-p03-060-fp-in-lomega} show that
\(\operatorname{lfp}(\Phi)=\bigcup_n\Phi^n\) and
\(\operatorname{gfp}(\Phi)=\bigcap_n\Phi_n\) on finite structures. The
displayed formulas use countable disjunctions and conjunctions,
\(\forall v_i\) and \(\neg\) over \(v_1,\ldots,v_k\) only, so they lie in
\(\mathrm{L}^k_{\infty\omega}(\sigma)\).
\end{proof}
\end{annexitem}

\begin{annexitem}{annex:p03-w03-p03-060-datalog-stages}{thm:p03-w03-p03-060-datalog-stages}{The stages of a Datalog program reach the least fixed point by stage \(\omega\)}
\begin{proof}
\emph{Monotonicity.} If \(\bar X\subseteq\bar Y\) componentwise and a body
\(\beta\) holds under \(\gamma\) in \((\mathbf A,\bar X)\), it holds in
\((\mathbf A,\bar Y)\): head atoms occur only unnegated and their relations
grew, while input literals, equalities and inequalities do not depend on
the head relations. Hence \(\Gamma(\bar X)\subseteq\Gamma(\bar Y)\).

\emph{The stages increase.} \(\bar X^0\subseteq\bar X^1\) trivially, and
\(\bar X^m\subseteq\bar X^{m+1}\) gives
\(\bar X^{m+1}=\Gamma(\bar X^m)\subseteq\Gamma(\bar X^{m+1})=\bar X^{m+2}\).

\emph{\(\bar X^\omega\) is a fixed point.} For every \(m\),
\(\bar X^m\subseteq\bar X^\omega\), so
\(\bar X^{m+1}=\Gamma(\bar X^m)\subseteq\Gamma(\bar X^\omega)\); with
\(\bar X^0=\varnothing\) this gives
\(\bar X^\omega\subseteq\Gamma(\bar X^\omega)\). Conversely, let
\(\bar a\in\Gamma_i(\bar X^\omega)\) by a rule and a map \(\gamma\). The
body has finitely many head atoms, and the value tuple of each lies in the
corresponding component of \(\bar X^\omega\), hence in that component of
some \(\bar X^{m_j}\). Let \(m\) be the largest \(m_j\) (or \(0\) if there
are none). Since the stages increase, all these tuples lie in
\(\bar X^m\), so the body holds under \(\gamma\) in
\((\mathbf A,\bar X^m)\) and \(\bar a\in\Gamma_i(\bar X^m)=X^{m+1}_i\subseteq X^\omega_i\).

\emph{Leastness.} If \(\Gamma(\bar Y)\subseteq\bar Y\), then
\(\bar X^0\subseteq\bar Y\), and \(\bar X^m\subseteq\bar Y\) gives
\(\bar X^{m+1}=\Gamma(\bar X^m)\subseteq\Gamma(\bar Y)\subseteq\bar Y\); so
\(\bar X^\omega\subseteq\bar Y\). Every fixed point satisfies this
hypothesis, so \(\bar X^\omega\) is the least fixed point.
\end{proof}
\end{annexitem}

\begin{annexitem}{annex:p03-w03-p03-060-datalog-in-existential}{thm:p03-w03-p03-060-datalog-in-existential}{Datalog is contained in existential finite-variable infinitary logic}
\begin{proof}
\emph{Normal form of the rules.} Let \(q\) be the largest number of
variables in a rule and \(w\) the largest head arity. Renaming the
variables of a single rule bijectively does not change \(\Gamma\), which
quantifies over all maps \(\gamma\); so we may assume that the variables of
every rule are among \(y_1,\ldots,y_q\), and we choose distinct variables
\(x_1,\ldots,x_w,z_1,\ldots,z_w\) different from them.

\emph{Update formulas.} For a rule \(\varrho=R_i(u_1,\ldots,u_{r_i})\leftarrow\beta\)
let
\[
 \delta_\varrho\defeq\exists y_1\cdots\exists y_q\bigl(\textstyle\bigwedge_{j\le r_i}x_j=u_j\wedge\bigwedge\beta\bigr),
\]
where \(\bigwedge\beta\) is the conjunction of the body literals, and let
\(\delta_i\defeq\bigvee\{\delta_\varrho:\varrho\text{ has head symbol }R_i\}\)
(\(\bot\) if there is none). For every \(\bar X\),
\(\Gamma_i(\bar X)=\{\bar a:(\mathbf A,\bar X)\models\delta_i[\bar x\mapsto\bar a]\}\):
a tuple \(\bar a\) satisfies \(\delta_\varrho\) exactly when some values of
\(y_1,\ldots,y_q\), that is some \(\gamma\), make the body true with
\(a_j=\gamma(u_j)\) (the \(x_j\) are not among the \(y\)'s, so the values
\(a_j\) are not disturbed). Each \(\delta_i\) is in
\(\exists\mathrm{L}_{\infty\omega}\) (with \(\neq\) for a
Datalog\((\neq)\) program), its free variables are among
\(x_1,\ldots,x_{r_i}\), and its variables are among the \(x\)'s and
\(y\)'s.

\emph{Stage formulas.} Put \(\varphi_i^0\defeq\bot\) and
\(\varphi_i^{m+1}\defeq\delta_i[R_1,\ldots,R_n:=\varphi_1^m,\ldots,\varphi_n^m]\).
By induction on \(m\), the free variables of \(\varphi_l^m\) are among
\(x_1,\ldots,x_{r_l}\), so there are no parameters, no \(z_j\) is free in
\(\varphi_l^m\), and no \(z_j\) occurs in \(\delta_i\);
\cref{thm:p03-w03-p03-060-substitution} applies and gives, for every
\(\sigma\)-structure \(\mathbf A\),
\[
 \{\bar a:\mathbf A\models\varphi_i^{m+1}[\bar x\mapsto\bar a]\}
 =\Gamma_i\bigl(\{\bar a:\mathbf A\models\varphi_l^m[\bar x\mapsto\bar a]\}_{l\le n}\bigr).
\]
Hence \(\varphi_i^m\) defines \(X^m_i\) for every \(m\). The variables of
every \(\varphi_i^m\) are among \(y_1,\ldots,y_q,x_1,\ldots,x_w,z_1,\ldots,z_w\),
and the lemma keeps \(\varphi_i^m\) in \(\exists\mathrm{L}_{\infty\omega}\)
(with \(\neq\)).

\emph{Conclusion.} Let \(\varphi_i\defeq\bigvee_{m\in\N}\varphi_i^m\). It
defines \(\bigcup_mX^m_i\), which is \((\Pi,R_i)(\mathbf A)\) by
\cref{def:p03-w03-p03-060-datalog} and, by
\cref{thm:p03-w03-p03-060-datalog-stages}, the \(i\)-th component of the
least fixed point of \(\Gamma\), on every structure. It has at most
\(q+2w\) variables.
\end{proof}
\end{annexitem}

\begin{annexitem}{annex:p03-w03-p03-060-datalog-in-lfp}{thm:p03-w03-p03-060-datalog-in-lfp}{Datalog is a fragment of LFP}
\begin{proof}
Let \(M=\mathcal P(U_1)\times\cdots\times\mathcal P(U_n)\) with finite
\(U_l\), ordered componentwise.

\emph{Lemma A.} If \(F\colon M\to M\) is monotone, its stages
\(\bar X^0=(\varnothing,\ldots,\varnothing)\),
\(\bar X^{m+1}=F(\bar X^m)\) increase and are constant from some \(m_0\)
on, and \(\bar X^{m_0}\) is the least fixed point of \(F\) and lies below
every \(\bar Y\) with \(F(\bar Y)\le\bar Y\). \emph{Proof.} The induction
for increase and the induction \(\bar X^m\le\bar Y\) are those in the proof
of \cref{thm:p03-w03-p03-060-datalog-stages}; an increasing sequence in the
finite set \(M\) is eventually constant, and its constant value
\(\bar X^{m_0}\) satisfies \(F(\bar X^{m_0})=\bar X^{m_0+1}=\bar X^{m_0}\).
Every fixed point satisfies \(F(\bar Y)\le\bar Y\). Write
\(\operatorname{lfp}(F)\) for \(\bar X^{m_0}\).

\emph{Lemma B (Beki\'c principle).} Let \(M_1=\mathcal P(U_1)\) and
\(M_2=\mathcal P(U_2)\times\cdots\times\mathcal P(U_n)\), and let
\(F\colon M_1\times M_2\to M_1\), \(G\colon M_1\times M_2\to M_2\) be
monotone. For \(X\in M_1\) let \(G_X(\bar Y)=G(X,\bar Y)\), monotone on
\(M_2\), and \(E(X)\defeq F(X,\operatorname{lfp}(G_X))\). Then \(E\) is
monotone and \(\operatorname{lfp}(E)\) is the first component \(X^\ast\)
of \(\operatorname{lfp}(F,G)=(X^\ast,\bar Y^\ast)\). \emph{Proof.} If
\(X\subseteq X'\), then
\(G_X(\operatorname{lfp}G_{X'})\le G_{X'}(\operatorname{lfp}G_{X'})=\operatorname{lfp}G_{X'}\),
so \(\operatorname{lfp}G_X\le\operatorname{lfp}G_{X'}\) by Lemma~A, and
\(E(X)\subseteq E(X')\) by monotonicity of \(F\). Next,
\(G_{X^\ast}(\bar Y^\ast)=\bar Y^\ast\), so
\(\operatorname{lfp}G_{X^\ast}\le\bar Y^\ast\) and
\(E(X^\ast)\subseteq F(X^\ast,\bar Y^\ast)=X^\ast\); by Lemma~A,
\(\operatorname{lfp}(E)\subseteq X^\ast\). Conversely, let
\(X_0=\operatorname{lfp}(E)\) and \(\bar Y_0=\operatorname{lfp}(G_{X_0})\).
Then \(F(X_0,\bar Y_0)=E(X_0)=X_0\) and
\(G(X_0,\bar Y_0)=G_{X_0}(\bar Y_0)=\bar Y_0\), so \((X_0,\bar Y_0)\) is a
fixed point of \((F,G)\) and \(X^\ast\subseteq X_0\) by Lemma~A.

\emph{Claim.} Let \(\tau\) be a vocabulary and
\(\delta_1(\bar x_1),\ldots,\delta_n(\bar x_n)\) first-order formulas of
vocabulary \(\tau\cup\{R_1,\ldots,R_n\}\), with
\(\bar x_l=(x_1,\ldots,x_{r_l})\) for fixed distinct \(x_1,\ldots,x_w\),
free variables of \(\delta_l\) among \(\bar x_l\), every \(R_j\) occurring
only positively, no universal quantifier, and negation applied only to
atoms of \(\tau\) and to equalities. For every \(i\le n\) there is an LFP
formula \(\chi_i\) of vocabulary \(\tau\) with free variables among
\(\bar x_i\), built with \(\mathbf{lfp}\) only, with the same restriction
on \(\forall\) and negation inside every fixed-point body and outside, in
which every \(P\in\tau\) that occurs only positively in all the
\(\delta_l\) occurs only positively, and which on every finite
\(\tau\)-structure \(\mathbf A\) defines the \(i\)-th component of
\(\operatorname{lfp}(\Delta)\), where
\(\Delta(\bar X)=(\{\bar a:(\mathbf A,\bar X)\models\delta_l[\bar x_l\mapsto\bar a]\})_{l\le n}\).
(\(\Delta\) is monotone: enlarging the interpretation of one \(R_j\) at a
time, each step enlarges every component by
\cref{thm:p03-w03-p03-059-positive-monotone}, since \(R_j\) is positive in
every \(\delta_l\).)

Induction on \(n\). For \(n=1\), \(\chi_1\defeq[\mathbf{lfp}\,R_1\bar x_1\,.\,\delta_1](\bar x_1)\)
defines \(\operatorname{lfp}(\Delta)\) by \cref{def:p03-w03-p03-059-lfp},
and the occurrences of each \(P\) are those of \(\delta_1\). For \(n>1\)
and given \(i\), renumber so that \(i=1\); renumbering the components
renumbers the least fixed point. Apply the induction hypothesis to the
system \(\delta_2,\ldots,\delta_n\) over the vocabulary
\(\tau\cup\{R_1\}\), in which \(R_1\) occurs only positively: it gives
\(\chi'_j\) (\(2\le j\le n\)) of vocabulary \(\tau\cup\{R_1\}\), with
\(R_1\) only positive, defining on each finite \((\mathbf A,X)\) the
\(j\)-th component of \(\operatorname{lfp}(G_X)\), where \(G\) is the
operator of \(\delta_2,\ldots,\delta_n\) and \(X\) interprets \(R_1\). Choose
new variables \(z_1,\ldots,z_w\) and let
\(\theta\defeq\delta_1[R_2,\ldots,R_n:=\chi'_2,\ldots,\chi'_n]\) as in
\cref{thm:p03-w03-p03-060-substitution}; there are no parameters. By that
lemma, \(R_1\) and every \(P\) positive in all \(\delta_l\) occur only
positively in \(\theta\), the restriction on \(\forall\) and negation is
kept (the replacement uses \(\exists\), \(\wedge\) and equality), and on
every finite \((\mathbf A,X)\)
\[
 \{\bar a:(\mathbf A,X)\models\theta[\bar x_1\mapsto\bar a]\}
 =F_1(X,\operatorname{lfp}(G_X))=E(X),
\]
where \(F_1\) is the operator of \(\delta_1\). Put
\(\chi_1\defeq[\mathbf{lfp}\,R_1\bar x_1\,.\,\theta](\bar x_1)\). It
defines \(\operatorname{lfp}(E)\), which by Lemma~B is the first component
of \(\operatorname{lfp}(\Delta)=\operatorname{lfp}(F_1,G)\).

\emph{Application.} For a Datalog\((\neq)\) program, possibly with negated
input atoms, take \(\tau=\sigma\) and the update formulas \(\delta_i\) of
the proof of \cref{thm:p03-w03-p03-060-datalog-in-existential} (that proof
did not use the absence of negated input atoms in showing
\(\Gamma_i(\bar X)=\{\bar a:(\mathbf A,\bar X)\models\delta_i[\bar x\mapsto\bar a]\}\)).
They are first-order, the head symbols occur only positively, they have no
universal quantifier, and negation is applied only to input atoms and
equalities. Then \(\Delta=\Gamma\), and by
\cref{thm:p03-w03-p03-060-datalog-stages} and Lemma~A the query
\((\Pi,R_i)\) is the \(i\)-th component of \(\operatorname{lfp}(\Gamma)\)
on every finite structure. The Claim gives \(\chi_i\).
\end{proof}
\end{annexitem}

\begin{annexitem}{annex:p03-w03-p03-060-continuity}{thm:p03-w03-p03-060-continuity}{Existential formulas are continuous, and their systems close at \(\omega\)}
\begin{proof}
\emph{Item 1.} By induction on \(\varphi\), for every assignment: if
\(\bar P\subseteq\bar P'\) componentwise and
\((\mathbf A,\bar P)\models\varphi[s]\), then
\((\mathbf A,\bar P')\models\varphi[s]\). An atom \(S_j(\bar t)\) holds
when the values of \(\bar t\) lie in \(P_j\subseteq P'_j\); other atoms and
inequalities do not mention \(\bar S\); conjunctions, disjunctions and
\(\exists\) preserve the implication.

\emph{Item 2.} Put \(\bar P^\omega\defeq\bigcup_n\bar P^n\). If
\((\mathbf A,\bar P^n)\models\varphi[s]\), then
\((\mathbf A,\bar P^\omega)\models\varphi[s]\) by item~1. The converse is
proved by induction on \(\varphi\), for all \(s\). An atom
\(S_j(\bar t)\) true in \((\mathbf A,\bar P^\omega)\) has its value tuple
in some \(P^n_j\). Atoms of \(\sigma\), equalities and inequalities do not
depend on \(\bar P\) (take \(n=0\)). If a disjunction is true, one
disjunct is true at \(\bar P^\omega\), hence at some \(\bar P^n\). If a
finite conjunction \(\bigwedge\Psi\) is true, each \(\psi\in\Psi\) is true
at some \(\bar P^{n_\psi}\); with \(n\) the largest \(n_\psi\) (or \(0\) if
\(\Psi=\emptyset\)), all are true at \(\bar P^n\) by item~1, as the
sequence increases. If \(\exists x\,\psi\) is true at \(s\), then \(\psi\)
is true at some \(s[x\mapsto a]\), hence at some \(\bar P^n\).

\emph{Item 3.} By item~1, \(\Phi\) is monotone. The stages increase:
\(\Phi^0\subseteq\Phi^1\), and \(\Phi^m\subseteq\Phi^{m+1}\) gives
\(\Phi^{m+1}=\Phi(\Phi^m)\subseteq\Phi(\Phi^{m+1})=\Phi^{m+2}\). Applying
item~2 to the chain \(\Phi^0\subseteq\Phi^1\subseteq\cdots\) gives
\(\Phi_i(\bigcup_m\Phi^m)=\bigcup_m\Phi_i(\Phi^m)=\bigcup_m\Phi^{m+1}_i=\bigcup_m\Phi^m_i\),
the last step because \(\Phi^0_i=\varnothing\). So \(\bigcup_m\Phi^m\) is a
fixed point. If \(\Phi(\bar Y)\subseteq\bar Y\), then
\(\Phi^0\subseteq\bar Y\), and \(\Phi^m\subseteq\bar Y\) gives
\(\Phi^{m+1}\subseteq\Phi(\bar Y)\subseteq\bar Y\); so
\(\bigcup_m\Phi^m\) lies below every fixed point and is the least one. No
finiteness of \(\mathbf A\) was used.
\end{proof}
\end{annexitem}

\begin{annexitem}{annex:p03-w03-p03-060-kv-systems}{thm:p03-w03-p03-060-kv-systems}{Kolaitis and Vardi: systems of positive \(\exists\mathrm{FO}^k\)-formulas}
\begin{proof}
For formulas \(\theta_1,\ldots,\theta_l\in\exists\mathrm{FO}^k(\sigma)\)
with \(\operatorname{fr}(\theta_j)\subseteq[n_j]\) and a formula \(\chi\)
of vocabulary \(\sigma\cup\{S_1,\ldots,S_l\}\) over \(v_1,\ldots,v_k\), let
\(\chi[\bar S:=\bar\theta]\) replace each atomic subformula
\(S_j(t_1,\ldots,t_{n_j})\) by \((\theta_j)_{\rho}\), where
\(\rho(c)\defeq t_c\) for \(c\le n_j\) and \(\rho(c)\defeq v_c\) otherwise
(\cref{thm:p03-w03-p03-060-variable-substitution}). Put
\(\theta_j^{\mathbf A}\defeq\{(\alpha(1),\ldots,\alpha(n_j)):\mathbf A\models\theta_j[\alpha]\}\).
The argument of the claim in the proof of
\cref{thm:p03-w03-p03-060-kv-stages}, with one symbol replaced by
several, gives for every \(\alpha\colon[k]\to A\)
\[
 \mathbf A\models\chi[\bar S:=\bar\theta][\alpha]\iff(\mathbf A,\theta_1^{\mathbf A},\ldots,\theta_l^{\mathbf A})\models\chi[\alpha],
\]
and \(\chi[\bar S:=\bar\theta]\) has no free label outside
\(\operatorname{fr}(\chi)\). If \(\chi\in\exists\mathrm{FO}^k\) (with
\(\neq\)), so is \(\chi[\bar S:=\bar\theta]\), because the lemma keeps each
\((\theta_j)_\rho\) in \(\exists\mathrm{FO}^k\) (with \(\neq\)) and the
replacement adds no constructor.

\emph{Item 1.} Put \(\varphi^0_i\defeq\bot\defeq\bigvee\emptyset\), which
belongs to \(\exists\mathrm{FO}^k\), and
\(\varphi^{m+1}_i\defeq\varphi_i[\bar S:=\varphi^m_1,\ldots,\varphi^m_l]\).
By induction on \(m\), \(\varphi^m_i\in\exists\mathrm{FO}^k(\sigma)\) has
free labels in \([n_i]\), and \((\varphi^m_i)^{\mathbf A}=\Phi^m_i\): for
\(m=0\) both are empty, and
\((\varphi^{m+1}_i)^{\mathbf A}=\{\bar a:(\mathbf A,\Phi^m)\models\varphi_i[\bar v\mapsto\bar a]\}=\Phi_i(\Phi^m)=\Phi^{m+1}_i\).

\emph{Item 2.} The formulas \(\varphi_i\) are built from atomic formulas
(and inequalities) by \(\exists\) and finite conjunctions and
disjunctions, so \cref{thm:p03-w03-p03-060-continuity} applies: \(\Phi\)
is monotone and \(\varphi^\infty_i=\bigcup_m\Phi^m_i\) on every structure.
The disjunction \(\bigvee_m\varphi^m_i\) is a formula of
\(\exists\mathrm{L}^k_{\infty\omega}(\sigma)\) (a countable disjunction of
\(\exists\mathrm{FO}^k\)-formulas over \(v_1,\ldots,v_k\)) with free labels
in \([n_i]\), and it holds at \(\alpha\) exactly when
\((\alpha(1),\ldots,\alpha(n_i))\in\bigcup_m\Phi^m_i\). The version with
\(\neq\) is the same proof.
\end{proof}
\end{annexitem}

\begin{annexitem}{annex:p03-w03-p03-060-kdatalog-system}{thm:p03-w03-p03-060-kdatalog-system}{\(k\)-Datalog programs as systems of positive \(\exists\mathrm{FO}^k\)-formulas}
\begin{proof}
For a tuple \(\bar X=(X_1,\ldots,X_n)\) of interpretations of the head
symbols let \(P(\bar X)\) be the tuple with components
\(P(\bar X)_{i,\mu}\defeq\{\alpha\in A^k:\mu^\alpha\in X_i\}\).

\emph{The formula of a rule.} Let \(\varrho\) be a rule
\(R_i(u_1,\ldots,u_{r_i})\leftarrow\beta_1\wedge\cdots\wedge\beta_q\), let
\(W\) be the set of variables of its body (\(|W|\le k\)), and fix an
injection \(\pi\colon W\to\{v_1,\ldots,v_k\}\). For a term \(t\) of the
body let \(t^\pi\defeq\pi(t)\) if \(t\) is a variable and \(t^\pi\defeq t\)
if it is a constant. Translate each body literal: an input atom, equality
or inequality by applying \(\pi\) to its terms; a head atom
\(R_l(t_1,\ldots,t_{r_l})\) by the atom \(S_{l,\nu}(v_1,\ldots,v_k)\) with
\(\nu(d)\defeq t_d^\pi\) (\(d\le r_l\)), a member of \(M_l\). Let
\(B_\varrho\) be the conjunction of the translated literals, and
\(d_1,\ldots,d_p\) the labels \(\pi(w)\) for the variables \(w\in W\) that
do not occur in the head. For \(\mu\in M_i\) let \(\hat\rho\colon[k]\to T_k\)
send the label \(\pi(w)\) of a head variable \(w\in W\) to \(\mu(j_w)\),
where \(j_w\) is the least \(j\) with \(u_j=w\), and every other label
\(d\) to \(v_d\). Put
\[
 E_{\varrho,\mu}\defeq\textstyle\bigwedge\{\mu(j)=\mu(j'):u_j=u_{j'}\text{ is a variable}\}\wedge\bigwedge\{\mu(j)=c:u_j=c\text{ is a constant}\},
\]
\[
 \psi_{\varrho,\mu}\defeq E_{\varrho,\mu}\wedge\bigl(\exists v_{d_1}\cdots\exists v_{d_p}\,B_\varrho\bigr)_{\hat\rho},
 \qquad
 \varphi_{i,\mu}\defeq\textstyle\bigvee\{\psi_{\varrho,\mu}:\varrho\text{ has head symbol }R_i\}.
\]
These are positive \(\exists\mathrm{FO}^k\)-formulas (with \(\neq\) when
\(\Pi\) has inequalities) of vocabulary \(\sigma\cup\bar S\) by
\cref{thm:p03-w03-p03-060-variable-substitution}, over
\(v_1,\ldots,v_k\), so the system is as required.

\emph{The operator.} We show: for every \(\bar X\), \(i\), \(\mu\) and
\(\alpha\colon[k]\to A\),
\((\mathbf A,P(\bar X))\models\varphi_{i,\mu}[\alpha]\iff\mu^\alpha\in\Gamma_i(\bar X)\).
Fix a rule \(\varrho\) with head symbol \(R_i\). For
\(\alpha'\colon[k]\to A\) let \(\gamma_{\alpha'}\colon W\to A\),
\(\gamma_{\alpha'}(w)\defeq\alpha'(\pi(w))\). Then \(B_\varrho\) holds at
\(\alpha'\) in \((\mathbf A,P(\bar X))\) iff the body of \(\varrho\) holds
under \(\gamma_{\alpha'}\) in \((\mathbf A,\bar X)\): for an input
literal, \(t^\pi\) has value \(\gamma_{\alpha'}(t)\) at \(\alpha'\); for a
head atom, \(\alpha'\in P(\bar X)_{l,\nu}\) iff
\(\nu^{\alpha'}=(\gamma_{\alpha'}(t_1),\ldots,\gamma_{\alpha'}(t_{r_l}))\in X_l\).
By \cref{thm:p03-w03-p03-060-variable-substitution},
\((\exists v_{d_1}\cdots\exists v_{d_p}B_\varrho)_{\hat\rho}\) holds at
\(\alpha\) iff \(\exists v_{d_1}\cdots\exists v_{d_p}B_\varrho\) holds at
\(\hat\rho^\alpha\), that is, iff \(B_\varrho\) holds at some \(\alpha'\)
that agrees with \(\hat\rho^\alpha\) outside \(\{d_1,\ldots,d_p\}\). As
\(\alpha'\) ranges over these, \(\gamma_{\alpha'}\) ranges over all maps
\(\gamma\colon W\to A\) with \(\gamma(w)=\mu(j_w)^{\mathbf A}[\alpha]\) for
the head variables \(w\in W\), because \(\pi\) is injective. Hence
\(\psi_{\varrho,\mu}\) holds at \(\alpha\) iff \(E_{\varrho,\mu}\) holds and
some such \(\gamma\) satisfies the body.

On the other hand \(\mu^\alpha\in\Gamma_i(\bar X)\) by the rule
\(\varrho\) iff there is a map \(\delta\) from the variables of
\(\varrho\) into \(A\) under which the body holds and
\(u_j^{\mathbf A}[\delta]=\mu(j)^{\mathbf A}[\alpha]\) for all \(j\). If
such \(\delta\) exists, \(E_{\varrho,\mu}\) holds (two head positions with
the same variable, or with a constant, force the corresponding values of
\(\mu\)), and \(\gamma\defeq\delta\restriction W\) is as required, since
\(\delta(w)=\mu(j_w)^{\mathbf A}[\alpha]\) for head variables in \(W\).
Conversely, given \(E_{\varrho,\mu}\) and \(\gamma\), let \(\delta\) extend
\(\gamma\) by \(\delta(h)\defeq\mu(j_h)^{\mathbf A}[\alpha]\) for the head
variables \(h\notin W\). For a head position \(j\) with a variable \(u_j=w\),
\(\delta(w)=\mu(j_w)^{\mathbf A}[\alpha]=\mu(j)^{\mathbf A}[\alpha]\) by
\(E_{\varrho,\mu}\); for a constant \(u_j=c\),
\(c^{\mathbf A}=\mu(j)^{\mathbf A}[\alpha]\) by \(E_{\varrho,\mu}\). The body
holds under \(\delta\), which agrees with \(\gamma\) on \(W\). Taking the
disjunction over the rules gives the assertion, that is,
\(\Phi(P(\bar X))=P(\Gamma(\bar X))\).

\emph{Stages and fixed point.} \(P(\varnothing,\ldots,\varnothing)\) is
empty in every component, so by induction
\(\Phi^m=P(\bar X^m)\) for every \(m\). By
\cref{thm:p03-w03-p03-060-continuity} the least fixed point of the system
is \(\bigcup_m\Phi^m=\bigcup_mP(\bar X^m)=P(\bigcup_m\bar X^m)\), the last
equality because \(\mu^\alpha\) lies in \(\bigcup_mX^m_i\) iff it lies in
some \(X^m_i\); and \(\bigcup_mX^m_i=(\Pi,R_i)(\mathbf A)\) by
\cref{def:p03-w03-p03-060-datalog}. Only the bound \(|W|\le k\) on the body
variables was used.
\end{proof}
\end{annexitem}

\begin{annexitem}{annex:p03-w03-p03-060-kdatalog}{thm:p03-w03-p03-060-kdatalog}{\(k\)-Datalog is contained in \(\exists\mathrm{L}^k_{\infty\omega}\), and \(k\)-Datalog\((\neq)\) in \(\exists\mathrm{L}^k_{\infty\omega}(\neq)\)}
\begin{proof}
Let \(\mu_0\in M_i\) be \(\mu_0(j)=v_j\) (\(j\le r_i\le k\)); then
\(\mu_0^\alpha=(\alpha(1),\ldots,\alpha(r_i))\). Apply
\cref{thm:p03-w03-p03-060-kv-systems} to the system of
\cref{thm:p03-w03-p03-060-kdatalog-system}. Its item~1 gives
\(\varphi^m_{i,\mu_0}\in\exists\mathrm{FO}^k\) with free labels in \([k]\)
defining \(\{\alpha:(\alpha(1),\ldots,\alpha(r_i))\in X^m_i\}\), and its
item~2 gives \(\chi'\in\exists\mathrm{L}^k_{\infty\omega}\) defining
\(\{\alpha:(\alpha(1),\ldots,\alpha(r_i))\in(\Pi,R_i)(\mathbf A)\}\). Let
\(\chi_i\defeq\exists v_{r_i+1}\cdots\exists v_k\,\chi'\) and define the
stage formulas likewise from \(\varphi^m_{i,\mu_0}\). Their free labels lie
in \([r_i]\). At \(\alpha\), \(\chi_i\) holds iff \(\chi'\) holds at some
\(\alpha'\) agreeing with \(\alpha\) on \([r_i]\), iff
\((\alpha(1),\ldots,\alpha(r_i))\in(\Pi,R_i)(\mathbf A)\), because the
condition depends only on the labels in \([r_i]\) and \(A\) is not empty.
The same holds for the stage formulas. The \((\neq)\) version uses the
\((\neq)\) versions of both results.
\end{proof}
\end{annexitem}

\begin{annexitem}{annex:p03-w03-p03-060-tc-not-fo}{thm:p03-w03-p03-060-tc-not-fo}{Least fixed points of \(\exists\mathrm{FO}^k\)-systems need infinitary disjunctions}
\begin{proof}
Suppose \(\theta(x,y)\) defines the transitive closure on all
\(\{E\}\)-structures. For \(n\ge1\) let
\(\pi_n(x,y)\defeq\exists z_0\cdots\exists z_n\bigl(z_0=x\wedge z_n=y\wedge\bigwedge_{1\le j\le n}E(z_{j-1},z_j)\bigr)\),
which says that there is a walk of length \(n\) from \(x\) to \(y\). Let
\(c,d\) be new constant symbols and
\(T\defeq\{\theta(c,d)\}\cup\{\neg\pi_n(c,d):n\ge1\}\). A finite subset of
\(T\) is contained in \(\{\theta(c,d)\}\cup\{\neg\pi_n(c,d):n\le N\}\) for
some \(N\). The directed path with vertices \(0,\ldots,N+1\), edges
\((j,j+1)\), \(c=0\) and \(d=N+1\) is a model: \(\theta(c,d)\) holds because
there is a path from \(0\) to \(N+1\), and every walk from \(0\) to
\(N+1\) has length exactly \(N+1\), since each edge raises the vertex by
one. By the compactness theorem (\cref{thm:p03-compactness-ultraproduct})
\(T\) has a model \((\mathbf G,c,d)\). There \(\theta(c,d)\) holds, so a
directed path of some length \(n\ge1\) leads from \(c\) to \(d\), and
\(\pi_n(c,d)\) holds, contradicting \(\neg\pi_n(c,d)\in T\).

The displayed formula belongs to \(\exists\mathrm{FO}^3\) and is positive
in \(S\); by \cref{thm:p03-w03-p03-060-continuity} its operator \(\Phi\) is
monotone and has a least fixed point \(L\). Let \(T\) be the transitive
closure. An edge, or an edge followed by a path, is a path, so
\(\Phi(T)\subseteq T\) and \(L\subseteq T\). Conversely, by induction on the
length \(n\ge1\) of a path \(a=a_0,a_1,\ldots,a_n=b\): for \(n=1\),
\((a,b)\in\Phi(L)=L\) by the first disjunct; for \(n+1\), the induction
hypothesis gives \((a_1,b)\in L\), and the second disjunct with witness
\(a_1\) gives \((a,b)\in\Phi(L)=L\). So \(L=T\), which by the first part is
defined by no first-order formula, in particular by no
\(\exists\mathrm{FO}^3\)-formula.
\end{proof}
\end{annexitem}

\begin{annexitem}{annex:p03-w03-p03-060-datalog-systems}{thm:p03-w03-p03-060-datalog-systems}{Datalog programs are existential-conjunctive systems}
\begin{proof}
\emph{Item 1.} Take the update formulas \(\delta_i(x_1,\ldots,x_{r_i})\)
of the proof of \cref{thm:p03-w03-p03-060-datalog-in-existential}. Each is
a finite disjunction, over the rules with head symbol \(R_i\), of
\(\exists y_1\cdots\exists y_q\bigl(\bigwedge_jx_j=u_j\wedge\bigwedge\beta\bigr)\),
whose matrix is a conjunction of equalities and of the body literals:
atoms, equalities, and inequalities if \(\Pi\) has them. Its free variables
are among the distinct variables \(x_1,\ldots,x_{r_i}\). That proof shows
\(\Gamma_i(\bar X)=\{\bar a:(\mathbf A,\bar X)\models\delta_i[\bar x\mapsto\bar a]\}\)
for every \(\bar X\), so the system \((\delta_1,\ldots,\delta_n)\) has the
operator \(\Gamma\). Equal operators have equal least fixed points, and
\((\Pi,R_i)(\mathbf A)\) is the \(i\)-th component of
\(\operatorname{lfp}(\Gamma)\) by \cref{thm:p03-w03-p03-060-datalog-stages}.

\emph{Item 2.} For each disjunct \(\exists z_1\cdots\exists z_p\,\psi\) of
\(\varphi_i\), choose distinct variables \(z'_1,\ldots,z'_p\) not occurring
in \(\varphi_i\), let \(\psi'\) be \(\psi\) with each \(z_c\) replaced by
\(z'_c\) (\(\psi\) has no quantifier, so every occurrence is replaced), and
form the rule \(S_i(\bar x_i)\leftarrow\psi'\), whose body literals are the
conjuncts of \(\psi'\). The free variables of the disjunct lie in
\(\bar x_i\), so the variables of the rule are those of \(\bar x_i\) and
\(z'_1,\ldots,z'_p\). A map \(\gamma\) on them makes the body true in
\((\mathbf A,\bar P)\) iff \(\psi\) holds at the assignment
\(\bar x_i\mapsto\gamma(\bar x_i)\), \(\bar z\mapsto\gamma(\bar z')\)
(the occurrences of each \(z_c\), and of each variable of \(\bar x_i\) not
among the \(z_c\), receive the same values). Hence the rule contributes to
\(\Gamma_i(\bar P)\) exactly the tuples \(\bar a\) with
\((\mathbf A,\bar P)\models\exists\bar z\,\psi[\bar x_i\mapsto\bar a]\), and
the rules for all disjuncts give \(\Gamma_i=\Phi_i\). Head symbols occur in
bodies only as atoms, so the program is a Datalog program, a
Datalog\((\neq)\) program if the system has inequalities. Then
\((\Pi,S_i)(\mathbf A)=\varphi^\infty_i\) as in item~1.

\emph{Item 3.} Without constant symbols all terms are variables. Take a
rule with an equality literal \(y=y'\). If \(y\) and \(y'\) are the same
variable, delete the literal. Otherwise replace \(y'\) by \(y\) everywhere
in the rule and delete the literal. The maps \(\gamma\) satisfying the old
body have \(\gamma(y)=\gamma(y')\), and \(\gamma\mapsto\gamma\restriction(\text{variables}\setminus\{y'\})\)
is a bijection onto the maps satisfying the new body, with inverse
extending by \(y'\mapsto\gamma(y)\); corresponding maps give the same head
value. So \(\Gamma\) is unchanged. Each step removes one literal; repeat.

\emph{Consequence.} If \(Q\) is defined on \(\mathcal C\) by \((\Pi,R_i)\),
item~1, after renumbering so that \(R_i\) is last, gives the system. If
\(Q\) is defined by \(\varphi^\infty_m\), item~2 gives the program
\((\Pi,S_m)\). Inequalities appear on one side exactly when they appear on
the other.
\end{proof}
\end{annexitem}

\begin{annexitem}{annex:p03-w03-p03-060-single-conjunction}{thm:p03-w03-p03-060-single-conjunction}{Single conjunctions do not suffice}
\begin{proof}
\emph{Item 1.} By item~2 of \cref{thm:p03-w03-p03-060-datalog-systems},
the system has the operator \(\Gamma\) of a Datalog\((\neq)\) program with
exactly one rule \(\varrho_i\) for each head symbol \(S_i\), and by
\cref{thm:p03-w03-p03-060-datalog-stages} its least fixed point is the
union of the increasing stages. Term values in the product are pairs:
\(t^{\mathbf A\times\mathbf B}[\langle\gamma_A,\gamma_B\rangle]=(t^{\mathbf A}[\gamma_A],t^{\mathbf B}[\gamma_B])\),
for variables by definition and for constants by the definition of
\(c^{\mathbf A\times\mathbf B}\). We show by induction on \(m\): if
\(\bar a\in\Phi^m_i(\mathbf A)\) and \(\bar b\in\Phi^m_i(\mathbf B)\), then
\(\langle\bar a,\bar b\rangle\in\Phi^m_i(\mathbf A\times\mathbf B)\). For
\(m=0\) there is nothing to show. If \(\bar a\in\Phi^{m+1}_i(\mathbf A)\),
some map \(\gamma_A\) on the variables of \(\varrho_i\) gives the head
value \(\bar a\) and makes the body true in \((\mathbf A,\Phi^m(\mathbf A))\);
likewise \(\gamma_B\) for \(\bar b\). Let
\(\gamma(w)\defeq(\gamma_A(w),\gamma_B(w))\). A \(\sigma\)-atom of the body
holds under \(\gamma\) because it holds in both coordinates; an equality
holds because both coordinates are equal; an inequality holds because the
first coordinates differ; an atom \(S_j(\bar t)\) holds by the induction
hypothesis. The head value under \(\gamma\) is
\(\langle\bar a,\bar b\rangle\). Finally, if
\(\bar a\in\varphi^\infty_i(\mathbf A)\) and
\(\bar b\in\varphi^\infty_i(\mathbf B)\), both lie in the stage with the
larger index, since the stages increase.

\emph{Item 2.} Let \(\mathbf G\) have universe \(\{a,b\}\) and the one edge
\((a,b)\), and \(\mathbf H\) universe \(\{c,d,e\}\) and edges \((c,d)\),
\((d,e)\). Then \((a,b)\) and \((c,e)\) are in the transitive closures.
An edge of \(\mathbf G\times\mathbf H\) from \((g,h)\) to \((g',h')\)
requires the edge \((g,g')\) of \(\mathbf G\), so \(g=a\) and \(g'=b\);
hence every directed path in \(\mathbf G\times\mathbf H\) has length
\(1\), and a path from \((a,c)\) to \((b,e)\) would be an edge, requiring
the edge \((c,e)\) in \(\mathbf H\), which is absent. So
\(\langle(a,b),(c,e)\rangle\) is not in the transitive closure of
\(\mathbf G\times\mathbf H\). By item~1 the transitive closure is
\(\varphi^\infty_i\) for no single-conjunction system on the finite graphs.
The two-rule program defines it: its operator \(\Gamma\) puts \((a,b)\)
into \(S\) exactly when \((a,b)\) is an edge or there is \(c\) with
\((a,c)\) an edge and \((c,b)\in S\), so \(\Gamma\) is the operator of the
formula \(E(v_1,v_2)\vee\exists v_3(E(v_1,v_3)\wedge S(v_3,v_2))\); its
query is \(\operatorname{lfp}\Gamma\) by
\cref{thm:p03-w03-p03-060-datalog-stages}, and that least fixed point is
the transitive closure by the argument in the proof of
\cref{thm:p03-w03-p03-060-tc-not-fo}.
\end{proof}
\end{annexitem}

\begin{annexitem}{annex:p03-w03-p03-060-efp-datalog}{thm:p03-w03-p03-060-efp-datalog}{Gr\"adel: Datalog is equivalent to EFP}
\begin{proof}
All lattices below are products \(\prod_{\iota\in I}\mathcal P(U_\iota)\),
ordered componentwise and identified with the power set of the disjoint
union of the \(U_\iota\), so that \cref{thm:p03-w03-p03-059-knaster-tarski}
applies: a monotone operator \(F\) has a least fixed point, which lies
below every \(Y\) with \(F(Y)\le Y\). The sets \(U_\iota\) and \(I\) may be
infinite.

\emph{Step 0: monotonicity.} By induction on \(\psi\in\mathrm{EFP}(\tau)\),
simultaneously with the semantics: if \(\psi\) has no literal
\(\neg P(\bar t)\), then \(\psi\) is monotone in \(P\). Literals, \(\wedge\),
\(\vee\) and \(\exists\) are as in item~1 of
\cref{thm:p03-w03-p03-060-continuity}. For
\([\mathbf{lfp}\,R_j:\mathcal S](\bar t)\): each \(\varphi_l\) has no
literal \(\neg R_n(\bar s)\), so by the induction hypothesis
\(F_{\mathcal S}\) is monotone in each \(R_n\), hence monotone, and the
least fixed point exists. If \(P\) is not negated in \(\psi\) and
\(\mathbf A'\) enlarges \(P^{\mathbf A}\), then
\(F^{\mathbf A}_{\mathcal S}(\bar Y)\le F^{\mathbf A'}_{\mathcal S}(\bar Y)\)
for all \(\bar Y\) by the induction hypothesis, so
\(F^{\mathbf A}_{\mathcal S}(\operatorname{lfp}F^{\mathbf A'}_{\mathcal S})\le\operatorname{lfp}F^{\mathbf A'}_{\mathcal S}\)
and \(\operatorname{lfp}F^{\mathbf A}_{\mathcal S}\le\operatorname{lfp}F^{\mathbf A'}_{\mathcal S}\).

\emph{Step 1: three lattice facts.} (L1) If
\(F((Y_\iota)_{\iota\in I})=(F_\iota(Y_\iota))_{\iota\in I}\) with each
\(F_\iota\) monotone, then
\(\operatorname{lfp}F=(\operatorname{lfp}F_\iota)_\iota\): this family is a
fixed point, and if \(F(Y)\le Y\) then \(F_\iota(Y_\iota)\le Y_\iota\), so
\(\operatorname{lfp}F_\iota\le Y_\iota\) for each \(\iota\). (L2) (Beki\'c.)
Let \(F\colon M_1\times M_2\to M_1\) and \(G\colon M_1\times M_2\to M_2\)
be monotone, \(G_X(Y)\defeq G(X,Y)\) and
\(E(X)\defeq F(X,\operatorname{lfp}G_X)\). Then \(E\) is monotone and
\(\operatorname{lfp}(F,G)=(X_0,\operatorname{lfp}G_{X_0})\) with
\(X_0\defeq\operatorname{lfp}E\). The proof of Lemma~B in
\cref{annex:p03-w03-p03-060-datalog-in-lfp} uses Lemma~A only through the
two properties of least fixed points recalled above, so it holds here and
shows that \(E\) is monotone, that \((X_0,Y_0)\) with
\(Y_0\defeq\operatorname{lfp}G_{X_0}\) is a fixed point of \((F,G)\), and
that the first component \(X^\ast\) of \(\operatorname{lfp}(F,G)\) equals
\(X_0\). The second component \(Y^\ast\) satisfies \(Y^\ast\le Y_0\) by
leastness and \(G_{X_0}(Y^\ast)=Y^\ast\), so \(Y_0\le Y^\ast\). (L3) If
\(F(Y,Z)=(F_1(Y),F_2(Y))\) with \(F_1,F_2\) monotone, then
\(\operatorname{lfp}F=(\operatorname{lfp}F_1,F_2(\operatorname{lfp}F_1))\):
this is (L2) with \(G(Y,Z)\defeq F_2(Y)\), for which
\(\operatorname{lfp}G_Y=F_2(Y)\).

For a program \(\Pi\) whose input symbols are interpreted in a structure
\(\mathbf B\), \(\Gamma_\Pi\) is monotone in the head relations
(\cref{thm:p03-w03-p03-060-datalog-stages}), and, by the same argument,
monotone in every input symbol that occurs in no negated body literal, and
\((\Pi,H)(\mathbf B)\) is the \(H\)-component of
\(\operatorname{lfp}\Gamma_\Pi\) on every structure. If two programs have
disjoint sets of head symbols and neither uses a head symbol of the other,
the operator of their union is the product of their operators, and by
(L1) the least fixed point of the union is the pair of least fixed points.
Adding rules with a new head symbol \(H\) that occurs in no body is the
situation of (L3): the old least fixed point is unchanged and the
\(H\)-component is the value of the new rules at it.

\emph{Item 1.} Let \(\mathcal S_\Pi\) be the system
\(R_l(x_1,\ldots,x_{r_l})\defeq\delta_l\) (\(l\le n\)) with the update
formulas \(\delta_l\) of the proof of
\cref{thm:p03-w03-p03-060-datalog-in-existential}. Each \(\delta_l\) is
built from the body literals and equalities by \(\wedge\), \(\vee\) and
\(\exists\); the head symbols occur in it unnegated; a negated input atom
or an inequality occurs in it only if it occurs in a rule. So
\(\delta_l\in\mathrm{EFP}(\sigma\cup\{R_1,\ldots,R_n\})\), with no
parameters, and
\([\mathbf{lfp}\,R_i:\mathcal S_\Pi](x_1,\ldots,x_{r_i})\in\mathrm{EFP}(\sigma)\).
That proof shows \(F_{\mathcal S_\Pi}=\Gamma\) on every structure (it does
not use the absence of negated input atoms); hence the formula defines the
\(i\)-th component of \(\operatorname{lfp}\Gamma\), which is
\((\Pi,R_i)(\mathbf A)\).

\emph{Item 2: lifting.} \emph{Lifting data} \(\mathfrak L\) are pairs
\((Q_1,\bar y_1),\ldots,(Q_r,\bar y_r)\) of distinct relation symbols
\(Q_e\in\tau\) and tuples \(\bar y_e\) of distinct variables. Let
\(\tau^{\mathfrak L}\) be \(\tau\) with each \(Q_e\) replaced by a new
symbol \(Q_e^+\) of arity \(\operatorname{ar}(Q_e)+|\bar y_e|\). For a
\(\tau^{\mathfrak L}\)-structure \(\mathbf B\) and an assignment \(s\)
defined on the \(\bar y_e\), let \(\mathbf B\langle s\rangle\) be the
\(\tau\)-structure with universe \(B\) that interprets \(Q_e\) as
\(\{\bar a:(\bar a,s(\bar y_e))\in(Q_e^+)^{\mathbf B}\}\) and the other
symbols as \(\mathbf B\). For a tuple \(\bar z\) and \(\bar c\) we write
\(\mathbf B\langle\bar c\rangle\) for \(\mathbf B\langle\bar z\mapsto\bar c\rangle\).

\emph{Claim.} Let \(\psi\in\mathrm{EFP}(\tau)\), let \(\mathfrak L\) be
lifting data no variable of which occurs in \(\psi\), let
\(\bar z=(z_1,\ldots,z_n)\) be distinct variables containing the free
variables of \(\psi\) and the variables of every \(\bar y_e\), and let
\(\mathcal N\) be a finite set of symbols. There are a Datalog program
\(\Pi\) with input symbols in \(\tau^{\mathfrak L}\) and head symbols
outside \(\tau^{\mathfrak L}\cup\mathcal N\), and a head symbol \(H\) of
arity \(n\), such that
\[
 \text{(E1)}\quad(\Pi,H)(\mathbf B)=\{\bar c\in B^n:\mathbf B\langle\bar c\rangle\models\psi[\bar z\mapsto\bar c]\}
 \quad\text{for every }\tau^{\mathfrak L}\text{-structure }\mathbf B,
\]
and (E2): a negated atom \(\neg P(\bar t)\) occurs in \(\Pi\) only if
\(\psi\) contains a negated \(P\)-literal (a negated \(Q_e\)-literal when
\(P=Q_e^+\)), and an inequality only if \(\psi\) contains a negated
equality. The proof is by induction on \(\psi\), for all
\(\mathfrak L\), \(\bar z\), \(\mathcal N\). Since a head symbol is never
an input symbol, programs built for different subformulas with disjoint
head symbols never use each other's heads.

\emph{Literal \(\lambda\).} Let \(\lambda^{\mathfrak L}\) be \(\lambda\) if
it mentions no \(Q_e\), and \(Q_e^+(\bar s,\bar y_e)\) or
\(\neg Q_e^+(\bar s,\bar y_e)\) if \(\lambda\) is \(Q_e(\bar s)\) or
\(\neg Q_e(\bar s)\). Let \(\Pi\defeq\{H(\bar z)\leftarrow\lambda^{\mathfrak L}\}\).
The variables of the rule are those of \(\bar z\). By (L3) with an empty
program, \((\Pi,H)(\mathbf B)\) is the set of \(\bar c\) at which
\(\lambda^{\mathfrak L}\) holds in \(\mathbf B\), which is the set at which
\(\lambda\) holds in \(\mathbf B\langle\bar c\rangle\), by the definition of
\(\mathbf B\langle\bar c\rangle\).

\emph{\(\psi_1\wedge\psi_2\), \(\psi_1\vee\psi_2\).} Apply the induction
hypothesis to \(\psi_1\) with \((\mathfrak L,\bar z,\mathcal N\cup\{H\})\)
and to \(\psi_2\) with
\((\mathfrak L,\bar z,\mathcal N\cup\{H\}\cup\text{heads of }\Pi_1)\). Let
\(\Pi\) be \(\Pi_1\cup\Pi_2\) with the rule
\(H(\bar z)\leftarrow H_1(\bar z)\wedge H_2(\bar z)\), respectively the two
rules \(H(\bar z)\leftarrow H_1(\bar z)\) and
\(H(\bar z)\leftarrow H_2(\bar z)\). By (L1) and (L3) the
\(H\)-component is the intersection, respectively the union, of the goal
components of \(\Pi_1\) and \(\Pi_2\).

\emph{\(\exists y\,\chi\).} The variable \(y\) occurs in \(\psi\), so it is
in no \(\bar y_e\). Let \(\bar z'\) be \(\bar z\) with \(y\) removed if it
occurs there; \(\bar z'\) still contains every \(\bar y_e\), and
\((\bar z',y)\) contains the free variables of \(\chi\). Apply the
induction hypothesis to \(\chi\) with \((\mathfrak L,(\bar z',y),\mathcal N\cup\{H\})\),
and add the rule \(H(\bar z)\leftarrow H_\chi(\bar z',w)\) for a variable
\(w\) not in \(\bar z\). Its variables are those of \(\bar z\) and \(w\).
By (L3) the \(H\)-component consists of the \(\bar c\) such that
\((\bar c\restriction\bar z',b)\) lies in the goal component of \(\Pi_\chi\)
for some \(b\), that is, such that
\(\mathbf B\langle\bar c\rangle\models\chi[\bar z'\mapsto\bar c\restriction\bar z',y\mapsto b]\)
for some \(b\) (the structure \(\mathbf B\langle\cdot\rangle\) depends only
on the values of the \(\bar y_e\), which lie in \(\bar z'\)). By
coincidence this is \(\mathbf B\langle\bar c\rangle\models\exists y\,\chi[\bar z\mapsto\bar c]\).

\emph{\([\mathbf{lfp}\,R_j:\mathcal S](\bar t)\).} Let \(\bar p\) list the
parameters of \(\mathcal S\) and, for each \(l\), \(\bar p_l\) those free in
\(\varphi_l\) outside \(\bar x_l\). Let \(\bar u\) list the distinct
variables of \(\bar y_1,\ldots,\bar y_r\), and let \(\bar p^\circ\) be a
tuple of \(|\bar p|\) distinct variables not occurring in \(\psi\) and not
in \(\bar z\); write \(p^\circ\) for the entry of \(\bar p^\circ\) at the
position of \(p\) in \(\bar p\). The lifting data
\(\mathfrak L_l\defeq\mathfrak L,(R_1,(\bar u,\bar p^\circ)),\ldots,(R_m,(\bar u,\bar p^\circ))\)
have no variable occurring in \(\varphi_l\); the tuple
\(\bar w_l\defeq(\bar x_l,\bar p_l,\bar u,\bar p^\circ)\) consists of
distinct variables and contains the free variables of \(\varphi_l\) and the
variables of all lifting tuples. The induction hypothesis for \(\varphi_l\),
of vocabulary \(\tau\cup\{R_1,\ldots,R_m\}\), with
\((\mathfrak L_l,\bar w_l,\mathcal N_l)\), where \(\mathcal N_l\) contains
\(\mathcal N\), \(H\), the \(R_n^+\) and the heads of
\(\Pi_1,\ldots,\Pi_{l-1}\), gives \(\Pi_l\) and \(G_l\) with input symbols
in \(\tau^{\mathfrak L}\cup\{R_1^+,\ldots,R_m^+\}\). By (E2) no \(R_n^+\)
occurs negated in \(\Pi_l\), because \(\varphi_l\) has no negated
\(R_n\)-literal. Let \(\Pi\) consist of the \(\Pi_l\), the rules
\[
 \varrho_l:\quad R_l^+(\bar x_l,\bar u,\bar p^\circ)\leftarrow G_l(\bar x_l,\bar p_l,\bar u,\bar p^\circ)\wedge\textstyle\bigwedge_{p\in\bar p_l}p=p^\circ\qquad(l\le m),
\]
in which \(R_1^+,\ldots,R_m^+\) become head symbols, and the rule
\(H(\bar z)\leftarrow R_j^+(\bar t,\bar u,\bar p)\), whose variables lie in
\(\bar z\) because the free variables of \(\psi\) are those of \(\bar t\)
and \(\bar p\). (E2) holds by the induction hypothesis.

Fix a \(\tau^{\mathfrak L}\)-structure \(\mathbf B\). The operator of
\(\Pi\) acts on triples \((\bar X,\bar Y,Z)\) (interpretations of the
\(R_l^+\), of the heads of the \(\Pi_l\), and of \(H\)) as
\((F(\bar Y),G(\bar X,\bar Y),K(\bar X))\), where \(G\) is the operator of
\(\bigcup_l\Pi_l\) on \((\mathbf B,\bar X)\), \(F_l\) is the operator of
\(\varrho_l\) and \(K\) that of the last rule; all are monotone. By (L3)
the \(H\)-component of \(\operatorname{lfp}\Gamma_\Pi\) is \(K(\bar X^\ast)\),
where \(\bar X^\ast\) is the first component of
\(\operatorname{lfp}(F,G)\), and by (L2) \(\bar X^\ast=\operatorname{lfp}E\)
with \(E(\bar X)=F(\operatorname{lfp}G_{\bar X})\). By (L1) and the
induction hypothesis (E1), the \(G_l\)-component of
\(\operatorname{lfp}G_{\bar X}\) is
\(\{\bar g:(\mathbf B,\bar X)\langle\bar g\rangle\models\varphi_l[\bar w_l\mapsto\bar g]\}\).
Write a tuple of \(R_l^+\) as \((\bar a,\bar d,\bar e)\) with \(\bar d\) for
\(\bar u\) and \(\bar e\) for \(\bar p^\circ\), and let
\(\bar X_{\bar d\bar e}\) be the family of slices
\(\{\bar a:(\bar a,\bar d,\bar e)\in X_n\}\). The equalities of
\(\varrho_l\) force the \(\bar p_l\)-entries to equal the corresponding
entries of \(\bar e\), and \((\mathbf B,\bar X)\langle\bar g\rangle\) is
\((\mathbf B\langle\bar d\rangle,\bar X_{\bar d\bar e})\). So
\[
 E(\bar X)_l=\{(\bar a,\bar d,\bar e):(\mathbf B\langle\bar d\rangle,\bar X_{\bar d\bar e})\models\varphi_l[\bar x_l\mapsto\bar a,\bar p\mapsto\bar e]\},
\]
using coincidence, since the free variables of \(\varphi_l\) are in
\(\bar x_l\) and \(\bar p_l\). Thus \(E\) acts slice by slice: its slice
at \((\bar d,\bar e)\) is the system operator \(F_{\mathcal S}\) computed in
\(\mathbf B\langle\bar d\rangle\) with the parameters \(\bar p\) assigned
\(\bar e\). By (L1), the slice of \(\operatorname{lfp}E\) at
\((\bar d,\bar e)\) is the least fixed point of that operator. Hence
\[
 K(\bar X^\ast)=\{\bar c:\bar t^{\mathbf B}[\bar c]\in\operatorname{lfp}\bigl(F_{\mathcal S}\text{ in }\mathbf B\langle\bar c\rangle,\ \bar p\mapsto\bar c\restriction\bar p\bigr)_j\}
 =\{\bar c:\mathbf B\langle\bar c\rangle\models[\mathbf{lfp}\,R_j:\mathcal S](\bar t)[\bar z\mapsto\bar c]\},
\]
since the values of \(\bar u\) under \(\bar z\mapsto\bar c\) determine
\(\mathbf B\langle\bar c\rangle\). This proves the Claim.

\emph{Item 2 and item 3.} Apply the Claim with empty lifting data, for
which \(\mathbf B\langle\bar c\rangle=\mathbf B\). Item~3 is (E2) for this
direction and the remark on \(\delta_l\) in item~1 for the other. The
three equivalences follow: every Datalog query is an EFP query by item~1,
every EFP query is a Datalog query by item~2, and item~3 matches the
kinds of negated literals on the two sides.
\end{proof}
\end{annexitem}

\begin{annexitem}{annex:p03-w03-p03-060-fp-comparison}{thm:p03-w03-p03-060-fp-comparison}{Least, inflationary and partial fixed points on finite structures}
\begin{proof}
Fix \(\mathbf A\) finite and the parameters, and let \(F=F_\psi\).

\emph{Item 1.} Let \(R^0=\varnothing\), \(R^{m+1}=F(R^m)\) and
\(S^0=\varnothing\), \(S^{m+1}=S^m\cup F(S^m)\). Since \(F\) is monotone,
the \(R^m\) increase and \(\bigcup_mR^m=\operatorname{lfp}(F)\) (proof of
\cref{thm:p03-w03-p03-060-fp-in-lomega}, case \(\mathbf{lfp}\)). By
induction \(S^m=R^m\): it holds for \(m=0\), and
\(S^{m+1}=R^m\cup F(R^m)=R^m\cup R^{m+1}=R^{m+1}\). So the inflationary
fixed point \(\bigcup_mS^m\) is \(\operatorname{lfp}(F)\), and the two
formulas agree at every \(\bar t\).

\emph{Item 2.} Let \(\psi'\defeq R(\bar x)\vee\psi\). Its operator is
\(F'(S)=S\cup F(S)\), because \(R(\bar x)\) holds at
\(s[\bar x\mapsto\bar a]\) in \((\mathbf A,S)\) exactly when
\(\bar a\in S\). The stages of \(F'\) from \(\varnothing\) are
\(T^{m+1}=T^m\cup F(T^m)\), the inflationary stages of \(\psi\). They
increase, so they are constant from some \(m_0\) on; then
\(F'(T^{m_0})=T^{m_0}\), a fixed point is reached, and the partial fixed
point of \(\psi'\) is \(T^{m_0}=\bigcup_mT^m\), the inflationary fixed
point of \(\psi\).

\emph{Translations.} Given an LFP formula, first replace each
\(\mathbf{gfp}\) subformula by the negated \(\mathbf{lfp}\) formula of
\cref{thm:p03-w03-p03-059-gfp-dual}. Then, by induction on formulas,
replace each \([\mathbf{lfp}\,R\bar x\,.\,\psi](\bar t)\) by
\([\mathbf{ifp}\,R\bar x\,.\,\psi^\dagger](\bar t)\), where \(\psi^\dagger\)
is the translation of \(\psi\). By the induction hypothesis \(\psi^\dagger\)
and \(\psi\) define the same operator on every finite structure, and it is
monotone by \cref{thm:p03-w03-p03-059-positive-monotone}, since \(R\) is
positive in the LFP formula \(\psi\); item~1 gives equivalence. For IFP to
PFP, replace each \([\mathbf{ifp}\,R\bar x\,.\,\psi](\bar t)\) by
\([\mathbf{pfp}\,R\bar x\,.\,R(\bar x)\vee\psi^\dagger](\bar t)\) and use
item~2; no monotonicity is needed.
\end{proof}
\end{annexitem}

\begin{annexitem}{annex:p03-w03-p03-060-stage-comparison}{thm:p03-w03-p03-060-stage-comparison}{Inflationary stage comparison is LFP-definable on finite structures}
\begin{proof}
Polarity of an occurrence (positive: under an even number of negations;
negative: under an odd number) is that of \cref{sec:p03-w03-p03-059-lfp},
counted through the bodies of fixed-point subformulas.

\emph{Step 0: renaming and replacement.}

(Relation binders.) Before making substitutions, rename every relation
variable bound by a fixed-point operator to a fresh symbol of the same
arity, using distinct symbols at distinct binders and avoiding \(X\) and
all free relation variables of \(\varphi\). At a binder
\([\mathbf{op}\,P\bar y\,.\,\theta](\bar t)\), rename only the occurrences
of \(P\) bound by this operator; a nested binder for \(P\) starts a new
scope. This preserves truth: for every interpretation \(Z\) of the bound
relation, a structural induction identifies the old body's interpretation
with that of the renamed body when the fresh symbol is interpreted by
\(Z\). The two bodies therefore define the same operator, and hence the
same least or greatest fixed point. Freshness prevents capture, and neither
the first-order parameters nor polarity changes. Consequently we may assume
that \(X\) is never bound inside \(\varphi\). All replacements of
\(X\)-atoms below concern occurrences free with respect to relation
binding; their first-order arguments may still be bound. This convention
also prevents the bound relation variables of a copied subformula from
capturing free relation variables at its insertion site.

(Renaming.) Let \(\xi\) be an LFP formula, \(y\) a variable and \(y'\) a
variable not occurring in \(\xi\); let \(\xi^{y\to y'}\) replace every
occurrence of \(y\), free or bound, including occurrences in the variable
tuples of fixed-point operators, by \(y'\). Then
\(\mathbf A\models\xi^{y\to y'}[s]\iff\mathbf A\models\xi[\tilde s]\) for
every structure (with every interpretation of the free relation variables)
and every assignment \(s\), where \(\tilde s\) sends \(y\) to \(s(y')\) and
agrees with \(s\) elsewhere. The induction is that of the Renaming step in
\cref{annex:p03-w03-p03-060-fp-in-lomega}, with one more case: for
\([\mathbf{op}\,P\bar y\,.\,\theta](\bar t)\), the renamed formula is
\([\mathbf{op}\,P\bar y^{y\to y'}\,.\,\theta^{y\to y'}](\bar t^{y\to y'})\);
by the induction hypothesis, for every interpretation \(Z\) of \(P\) and
every \(\bar a\), \(\theta^{y\to y'}\) holds at
\(s[\bar y^{y\to y'}\mapsto\bar a]\) exactly when \(\theta\) holds at
\(\tilde s[\bar y\mapsto\bar a]\) (the two assignments correspond under the
renaming), so the two bodies define the same operator, hence the same least
or greatest fixed point, and the terms have the same values. Consequently,
replacing a subformula \(Qy\,\xi\) (\(Q\in\{\exists,\forall\}\)) by
\(Qy'\,\xi^{y\to y'}\), or a subformula
\([\mathbf{op}\,P\bar y\,.\,\theta](\bar t)\) with \(y\) in \(\bar y\) by
\([\mathbf{op}\,P\bar y^{y\to y'}\,.\,\theta^{y\to y'}](\bar t)\), with
\(y'\) new, changes neither truth values nor free variables. Applying this
to every binding of a variable of \(\bar p\) inside \(\varphi\), we assume
from now on that \emph{no variable of \(\bar p\) is bound in \(\varphi\)};
the ranks, the stages and the relations to be defined do not change.

(Replacement.) Let \(\chi\) be an LFP formula and \(V\) a set of
first-order variables none of which is bound in \(\chi\). Let \(\chi'\)
arise from \(\chi\) by replacing some occurrences of atoms \(Q(\bar t)\) of
a relation symbol \(Q\) not bound in \(\chi\), by LFP formulas
\(\vartheta\) in which no relation
variable bound in \(\chi\) occurs free. Fix a structure \(\mathbf A\) and
an assignment \(s\), and suppose that every replaced atom \(Q(\bar t)\) and
its replacement \(\vartheta\) have the same truth value in \(\mathbf A\) at
every assignment that agrees with \(s\) on \(V\). Then
\(\mathbf A\models\chi[s]\iff\mathbf A\models\chi'[s]\). Proof: by
induction on subformulas \(\chi_1\) of \(\chi\) we show that \(\chi_1\) and
its image \(\chi_1'\) have the same truth value in every expansion
\(\mathbf A^+\) of \(\mathbf A\) by interpretations of the relation
variables bound in \(\chi\), at every assignment agreeing with \(s\) on
\(V\). A replaced atom: neither it nor \(\vartheta\) contains free a
relation variable bound in \(\chi\), so its truth value in \(\mathbf A^+\)
is that in \(\mathbf A\), and the hypothesis applies. Other atoms are
unchanged. \(\neg\), \(\wedge\), \(\vee\) commute with the replacement.
For \(\exists y\) and \(\forall y\), \(y\) is bound in \(\chi\), so
\(y\notin V\) and every \(s'[y\mapsto a]\) still agrees with \(s\) on
\(V\). For \([\mathbf{op}\,P\bar y\,.\,\theta](\bar t)\), the variables of
\(\bar y\) are bound in \(\chi\), so for every interpretation \(Z\) of
\(P\) and every \(\bar a\) the induction hypothesis applies to \(\theta\)
in \((\mathbf A^+,Z)\) at \(s'[\bar y\mapsto\bar a]\): \(\theta\) and
\(\theta'\) define the same operator, hence the same least or greatest
fixed point, and the terms \(\bar t\) are unchanged. Moreover \(\chi'\) is
an LFP formula: every relation variable bound by a fixed-point operator of
\(\chi\) still occurs only positively in its body, because the
\(\vartheta\) do not contain it free.

\emph{Step 1: ranks.} Fix a finite \(\tau\)-structure \(\mathbf A\) (with
interpretations of the free relation variables other than \(X\)) and an
assignment \(s\) to \(\bar p\). Let
\(F(S)\defeq\{\bar a\in A^r:(\mathbf A,S)\models\varphi[s[\bar x\mapsto\bar a]]\}\),
\(S^0\defeq\varnothing\), \(S^{m+1}\defeq S^m\cup F(S^m)\),
\(S^\infty\defeq\bigcup_mS^m\), and write \(|\bar a|\) for
\(|\bar a|_\varphi\).
\begin{enumerate}
\item[(R1)] The stages increase. If \(S^j=S^{j+1}\), then
      \(S^{j+2}=S^{j+1}\cup F(S^{j+1})=S^j\cup F(S^j)=S^{j+1}\), and by
      induction \(S^{j'}=S^j\) for all \(j'\ge j\). Hence, if
      \(|\bar a|=n+1\) with \(n\ge1\), then \(S^n\neq S^{n+1}\), so
      \(S^{n-1}\neq S^n\), and any \(\bar d\in S^n\setminus S^{n-1}\) has
      \(|\bar d|=n\).
\item[(R2)] If \(|\bar d|=n<\infty\), then
      \(S^n=\{\bar c:|\bar c|\le n\}\) and
      \(A^r\setminus S^n=\{\bar c:|\bar c|>n\}\), where \(\infty>n\).
\item[(R3)] \(\bar a\in S^{n+1}\) iff \(\bar a\in S^n\) or
      \(\bar a\in F(S^n)\).
\end{enumerate}
The target relations are
\(\lambda\defeq\{(\bar a,\bar b):|\bar a|\le|\bar b|,\ |\bar a|<\infty\}\)
and \(\eta\defeq\{(\bar a,\bar b):|\bar b|<|\bar a|\}\); in \(\eta\) the
rank of \(\bar b\) is finite. For a \(2r\)-ary relation \(T\) and
\(\bar d\in A^r\) write \(T_{\bar d}\defeq\{\bar c:(\bar c,\bar d)\in T\}\).
By (R2), if \(|\bar d|=n<\infty\) then \(\lambda_{\bar d}=S^n\) and
\(\eta_{\bar d}=A^r\setminus S^n\). Also \((\bar d,\bar d)\in\lambda\) iff
\(|\bar d|<\infty\).

\emph{Step 2: the update formulas.} Let \(L,H\) be new \(2r\)-ary relation
symbols, and \(\bar z=(z_1,\ldots,z_r)\) and \(y\) new variables (distinct
from each other, from \(\bar u,\bar w\), from \(\bar x\) and from the
variables of \(\varphi\)); let \(\bot\defeq\exists y\,\neg\,y=y\), a sentence false in
every structure.
\begin{itemize}
\item \(\varphi_0\) is \(\varphi\) with every free occurrence of \(X(\bar t)\) replaced by
      \(\bot\).
\item For \(\psi\in\{\varphi,\neg\varphi\}\), \(\psi^{\bar z}\) is \(\psi\)
      with every free occurrence of \(X(\bar t)\) at a positive position of \(\psi\)
      replaced by \(L(\bar t,\bar z)\) and every atom \(X(\bar t)\) at a
      negative position replaced by \(\neg H(\bar t,\bar z)\).
\item For an \(r\)-tuple \(\bar e\) of variables (here \(\bar u\) or
      \(\bar w\)) and a formula \(\xi\),
      \(\xi\langle\bar e\rangle\defeq\exists x_1\cdots\exists x_r\bigl(\bigwedge_{j\le r}x_j=e_j\wedge\xi\bigr)\).
\end{itemize}
Put
\begin{align*}
 \delta_L(\bar u,\bar w)&\defeq\varphi_0\langle\bar u\rangle\vee
 \exists\bar z\Bigl(L(\bar z,\bar z)\wedge\bigl(L(\bar u,\bar z)\vee\varphi^{\bar z}\langle\bar u\rangle\bigr)\wedge H(\bar w,\bar z)\Bigr),\\
 \delta_H(\bar u,\bar w)&\defeq\bigl(\varphi_0\langle\bar w\rangle\wedge\neg\varphi_0\langle\bar u\rangle\bigr)\vee{}\\
 &\phantom{{}\defeq{}}\exists\bar z\Bigl(L(\bar z,\bar z)\wedge\bigl(L(\bar w,\bar z)\vee\varphi^{\bar z}\langle\bar w\rangle\bigr)\wedge H(\bar u,\bar z)\wedge(\neg\varphi)^{\bar z}\langle\bar u\rangle\Bigr).
\end{align*}
Their free first-order variables are among \(\bar u,\bar w,\bar p\), and no
variable of \(\bar p\) is bound in them (the bound variables are those
bound in \(\varphi\), those of \(\bar x\), those of \(\bar z\), and \(y\)).

\emph{Positivity.} In \(\psi^{\bar z}\) each \(L\) stands at the position
of a positive occurrence of \(X\), and each \(H\) under one more negation
than a negative occurrence of \(X\); so \(L\) and \(H\) occur only
positively in \(\psi^{\bar z}\), and hence in \(\psi^{\bar z}\langle\bar e\rangle\).
\(\varphi_0\) contains neither symbol. So \(L\) and \(H\) occur only
positively in \(\delta_L\) and \(\delta_H\). Let
\(\Theta_L(L,H)\) be the set of \((\bar a,\bar b)\in A^{2r}\) with
\((\mathbf A,L,H)\models\delta_L[s[\bar u\mapsto\bar a,\bar w\mapsto\bar b]]\),
\(\Theta_H\) likewise with \(\delta_H\), and
\(\Theta(L,H)\defeq(\Theta_L(L,H),\Theta_H(L,H))\) on
\(\mathcal P(A^{2r})\times\mathcal P(A^{2r})\), ordered componentwise. By
\cref{thm:p03-w03-p03-059-positive-monotone}, applied to \(L\) with \(H\)
fixed and to \(H\) with \(L\) fixed, both components are monotone in each
argument, hence \(\Theta\) is monotone.

\emph{Step 3: reading the update formulas.} First, for an \(r\)-tuple
\(\bar e\) of variables not among \(\bar x\), \(\xi\langle\bar e\rangle\)
holds at an assignment \(s'\) iff \(\xi\) holds at
\(s'[\bar x\mapsto s'(\bar e)]\). Let \(\bar a,\bar d\in A^r\).
\begin{enumerate}
\item[(a)] \(\varphi_0\langle\bar u\rangle\) holds at
      \(\bar u\mapsto\bar a\) iff \(\bar a\in F(\varnothing)\): by
      Replacement with \(V=\emptyset\), \(\varphi_0\) holds in
      \(\mathbf A\) iff \(\varphi\) holds in \((\mathbf A,\varnothing)\),
      because an atom \(X(\bar t)\) is false when \(X\) is empty, like
      \(\bot\). The same holds with \(\bar w\) in place of \(\bar u\).
\item[(b)] Let \((L,H)\) satisfy \(L_{\bar d}=S\) and
      \(H_{\bar d}=A^r\setminus S\) for some \(S\subseteq A^r\). For
      \(\psi\in\{\varphi,\neg\varphi\}\), \(\psi^{\bar z}\langle\bar u\rangle\)
      holds in \((\mathbf A,L,H)\) at \(\bar u\mapsto\bar a\),
      \(\bar z\mapsto\bar d\) (and \(\bar p\) at its values under \(s\))
      iff \(\psi\) holds in \((\mathbf A,S)\) at \(s[\bar x\mapsto\bar a]\).
      In particular \(\varphi^{\bar z}\langle\bar u\rangle\) holds iff
      \(\bar a\in F(S)\), and \((\neg\varphi)^{\bar z}\langle\bar u\rangle\)
      holds iff \(\bar a\notin F(S)\). Indeed, apply Replacement to
      \(\psi\), read in \((\mathbf A,S,L,H)\), with
      \(V=\{z_1,\ldots,z_r\}\), which is not bound in \(\psi\): at every
      assignment sending \(\bar z\) to \(\bar d\), \(L(\bar t,\bar z)\)
      holds iff the value of \(\bar t\) lies in \(L_{\bar d}=S\), that is
      iff \(X(\bar t)\) holds, and \(\neg H(\bar t,\bar z)\) holds iff that
      value lies outside \(H_{\bar d}\), that is in \(S\). The free
      variables of \(\psi\) other than \(\bar x\) are among \(\bar p\). The
      same holds with \(\bar w\) in place of \(\bar u\).
\end{enumerate}

\emph{Step 4: \(\Theta(\lambda,\eta)\subseteq(\lambda,\eta)\).} Let
\((\bar a,\bar b)\in\Theta_L(\lambda,\eta)\). If the first disjunct of
\(\delta_L\) holds, \(\bar a\in F(\varnothing)=S^1\) by (a), so
\(|\bar a|=1\le|\bar b|\) and \((\bar a,\bar b)\in\lambda\). Otherwise
there is \(\bar d\) with \((\bar d,\bar d)\in\lambda\), so
\(n\defeq|\bar d|<\infty\), \(\lambda_{\bar d}=S^n\) and
\(\eta_{\bar d}=A^r\setminus S^n\); the middle conjunct says
\(\bar a\in S^n\) or, by (b), \(\bar a\in F(S^n)\), so
\(\bar a\in S^{n+1}\) by (R3) and \(|\bar a|\le n+1\); and
\((\bar b,\bar d)\in\eta\) says \(|\bar b|>n\), so
\(|\bar b|\ge n+1\ge|\bar a|\). Hence \((\bar a,\bar b)\in\lambda\).

Let \((\bar a,\bar b)\in\Theta_H(\lambda,\eta)\). If the first disjunct of
\(\delta_H\) holds, then by (a) \(\bar b\in S^1\) and \(\bar a\notin S^1\),
so \(|\bar b|=1<|\bar a|\). Otherwise there is \(\bar d\) with
\(n\defeq|\bar d|<\infty\); as before \(\bar b\in S^{n+1}\);
\((\bar a,\bar d)\in\eta\) gives \(\bar a\notin S^n\), and the last
conjunct gives \(\bar a\notin F(S^n)\) by (b); by (R3)
\(\bar a\notin S^{n+1}\), so \(|\bar a|>n+1\ge|\bar b|\) and
\((\bar a,\bar b)\in\eta\).

Let \((L^\ast,H^\ast)\) be the least fixed point of \(\Theta\) (Lemma~A in
\cref{annex:p03-w03-p03-060-datalog-in-lfp}, for the finite lattice
\(\mathcal P(A^{2r})^2\)). Since \((\lambda,\eta)\) satisfies
\(\Theta(\lambda,\eta)\subseteq(\lambda,\eta)\), that lemma gives
\(L^\ast\subseteq\lambda\) and \(H^\ast\subseteq\eta\).

\emph{Step 5: \(\lambda\subseteq L^\ast\) and \(\eta\subseteq H^\ast\).} We
prove by induction on \(m\ge1\) the statement \(P(m)\): every
\((\bar a,\bar b)\in\lambda\) with \(|\bar a|=m\) lies in \(L^\ast\), and
every \((\bar a,\bar b)\in\eta\) with \(|\bar b|=m\) lies in \(H^\ast\).
Every pair of \(\lambda\) has a first component of finite rank and every
pair of \(\eta\) a second component of finite rank, so this suffices.

First a consequence of the induction hypothesis. Suppose \(P(j)\) holds for
all \(j\le n\) and \(|\bar d|=n\). A pair \((\bar c,\bar d)\in\lambda\) has
\(|\bar c|\le n\), so it lies in \(L^\ast\) by \(P(|\bar c|)\); a pair
\((\bar c,\bar d)\in\eta\) has second component of rank \(n\), so it lies
in \(H^\ast\) by \(P(n)\). With Step~4,
\(L^\ast_{\bar d}=\lambda_{\bar d}=S^n\) and
\(H^\ast_{\bar d}=\eta_{\bar d}=A^r\setminus S^n\), so (b) applies to
\((L^\ast,H^\ast)\) and \(\bar d\) with \(S=S^n\). Also
\((\bar d,\bar d)\in\lambda\) has first component of rank \(n\), so
\((\bar d,\bar d)\in L^\ast\).

\(m=1\). If \((\bar a,\bar b)\in\lambda\) and \(|\bar a|=1\), then
\(\bar a\in F(\varnothing)\), the first disjunct of \(\delta_L\) holds by
(a), and \((\bar a,\bar b)\in\Theta_L(L^\ast,H^\ast)=L^\ast\). If
\((\bar a,\bar b)\in\eta\) and \(|\bar b|=1\), then
\(\bar b\in F(\varnothing)\) and \(|\bar a|>1\), so
\(\bar a\notin S^1=F(\varnothing)\); the first disjunct of \(\delta_H\)
holds, and \((\bar a,\bar b)\in\Theta_H(L^\ast,H^\ast)=H^\ast\).

\(m=n+1\) with \(n\ge1\), assuming \(P(1),\ldots,P(n)\).
If no tuple has rank \(n+1\), both assertions of \(P(n+1)\)
are vacuous. Otherwise (R1) supplies a tuple \(\bar d\) with
\(|\bar d|=n\), and the consequence above applies. Let
\((\bar a,\bar b)\in\lambda\) with \(|\bar a|=n+1\). Then
\((\bar d,\bar d)\in L^\ast\); \(\bar a\in S^{n+1}\setminus S^n\), so
\(\bar a\in F(S^n)\) by (R3) and \(\varphi^{\bar z}\langle\bar u\rangle\)
holds at \(\bar u\mapsto\bar a\), \(\bar z\mapsto\bar d\) in
\((\mathbf A,L^\ast,H^\ast)\) by (b); and \(|\bar b|\ge n+1>n=|\bar d|\),
so \((\bar b,\bar d)\in\eta\) with second component of rank \(n\), hence
\((\bar b,\bar d)\in H^\ast\) by \(P(n)\). The second disjunct of
\(\delta_L\) holds with witness \(\bar d\), and
\((\bar a,\bar b)\in L^\ast\). Let \((\bar a,\bar b)\in\eta\) with
\(|\bar b|=n+1\), so \(|\bar a|>n+1\). Then \((\bar d,\bar d)\in L^\ast\);
\(\bar b\in F(S^n)\), so \(\varphi^{\bar z}\langle\bar w\rangle\) holds at
\(\bar w\mapsto\bar b\), \(\bar z\mapsto\bar d\); \(|\bar a|>n\), so
\((\bar a,\bar d)\in\eta\) and \((\bar a,\bar d)\in H^\ast\) by \(P(n)\);
and \(\bar a\notin S^{n+1}\), so \(\bar a\notin F(S^n)\) and
\((\neg\varphi)^{\bar z}\langle\bar u\rangle\) holds by (b). The second
disjunct of \(\delta_H\) holds, and \((\bar a,\bar b)\in H^\ast\).

Thus \((L^\ast,H^\ast)=(\lambda,\eta)\).

\emph{Step 6: single fixed points.} Write
\(\gamma_H(\bar e,\bar e')\) for
\([\mathbf{lfp}\,H\bar u\bar w\,.\,\delta_H](\bar e,\bar e')\), where
\(\bar e,\bar e'\) are \(r\)-tuples of terms. \(L\) occurs in it only
positively, and its free first-order variables are those of
\(\bar e,\bar e'\) together with parameters among \(\bar p\). Let
\(\delta_L^\ast\) be \(\delta_L\) with every atom \(H(\bar e,\bar e')\)
replaced by \(\gamma_H(\bar e,\bar e')\). Every occurrence of \(H\) in
\(\delta_L\) is positive and \(L\) is positive in \(\gamma_H\), so \(L\)
occurs only positively in \(\delta_L^\ast\), and
\[
 \lambda_\varphi(\bar u,\bar w)\defeq[\mathbf{lfp}\,L\bar u\bar w\,.\,\delta_L^\ast](\bar u,\bar w)
\]
is an LFP formula with free first-order variables among
\(\bar u,\bar w,\bar p\). For a fixed interpretation \(L\), let
\(G_L(H)\defeq\Theta_H(L,H)\), a monotone operator on
\(\mathcal P(A^{2r})\); at every assignment agreeing with \(s\) on
\(\bar p\), \(\gamma_H(\bar e,\bar e')\) holds in \((\mathbf A,L)\) iff the
value of \((\bar e,\bar e')\) lies in \(\operatorname{lfp}(G_L)\).
Replacement with \(V\) the set of variables of \(\bar p\) (not bound in
\(\delta_L\)), applied in \((\mathbf A,L,\operatorname{lfp}(G_L))\), shows
that the operator of \(\delta_L^\ast\) is
\(E(L)\defeq\Theta_L(L,\operatorname{lfp}(G_L))\). By Lemma~B in
\cref{annex:p03-w03-p03-060-datalog-in-lfp} (with \(n=2\),
\(F=\Theta_L\), \(G=\Theta_H\)), \(\operatorname{lfp}(E)\) is the first
component \(L^\ast=\lambda\) of \(\operatorname{lfp}(\Theta)\). So
\(\lambda_\varphi\) defines \(\lambda\).

Exchanging the roles of \(L\) and \(H\) (Lemma~B applies to either
component, since the order on \(\mathcal P(A^{2r})^2\) is componentwise and
the two factors may be listed in either order): let
\(\gamma_L(\bar e,\bar e')\) be
\([\mathbf{lfp}\,L\bar u\bar w\,.\,\delta_L](\bar e,\bar e')\), in which
\(H\) is positive, let \(\delta_H^\ast\) be \(\delta_H\) with every atom
\(L(\bar e,\bar e')\) replaced by \(\gamma_L(\bar e,\bar e')\), and put
\(\eta_\varphi(\bar u,\bar w)\defeq[\mathbf{lfp}\,H\bar u\bar w\,.\,\delta_H^\ast](\bar u,\bar w)\).
The same argument shows that \(\eta_\varphi\) defines \(H^\ast=\eta\).

\emph{Step 7: the stage comparison relations.} By the definitions,
\(\bar a\le^{\mathrm{inf}}_\varphi\bar b\) iff \((\bar a,\bar b)\in\lambda\)
and \((\bar b,\bar b)\in\lambda\), and
\(\bar a\prec^{\mathrm{inf}}_\varphi\bar b\) iff
\((\bar b,\bar a)\in\eta\). Put
\(\mathrm{le}_\varphi(\bar u,\bar w)\defeq\lambda_\varphi(\bar u,\bar w)\wedge[\mathbf{lfp}\,L\bar u\bar w\,.\,\delta_L^\ast](\bar w,\bar w)\)
and
\(\mathrm{lt}_\varphi(\bar u,\bar w)\defeq[\mathbf{lfp}\,H\bar u\bar w\,.\,\delta_H^\ast](\bar w,\bar u)\).
These are LFP formulas with free first-order variables among
\(\bar u,\bar w,\bar p\), and they define \(\le^{\mathrm{inf}}_\varphi\) and
\(\prec^{\mathrm{inf}}_\varphi\); since \(\mathbf A\) and \(s\) were
arbitrary, this holds on every finite \(\tau\)-structure and for every
assignment to \(\bar p\).
\end{proof}
\end{annexitem}

\begin{annexitem}{annex:p03-w03-p03-060-stage-comparison-positive}{thm:p03-w03-p03-060-stage-comparison-positive}{Inflationary stage comparison is definable in positive LFP on finite structures}
\begin{proof}
\emph{Removing constants.} Let \(c_1,\ldots,c_s\) be the constant symbols
of \(\tau\) and \(\tau_0\) the vocabulary of its relation symbols (with the
free relation variables). Choose distinct variables \(q_1,\ldots,q_s\)
occurring nowhere in \(\varphi\) and distinct from
\(\bar u,\bar w,\bar x,\bar p\), and let \(\varphi'\) be \(\varphi\) with
every occurrence of each \(c_e\) replaced by \(q_e\). By induction on LFP
formulas \(\xi\): \((\mathbf A,S)\models\xi[s]\) iff
\((\mathbf A{\restriction}\tau_0,S)\models\xi'[s[\bar q\mapsto\bar c^{\mathbf A}]]\)
for every \(\tau\)-structure \(\mathbf A\), interpretation \(S\) of the
free relation variables and assignment \(s\). Term values agree; the
connectives and quantifiers commute (no \(q_e\) is bound, so updating a
bound variable keeps \(\bar q\mapsto\bar c^{\mathbf A}\)); and for a
fixed-point subformula the bodies define the same operator, hence the same
fixed point, while the terms \(\bar t\) and \(\bar t'\) have the same
values. So \(\varphi'\), an LFP formula of the relational vocabulary
\(\tau_0\cup\{X\}\) with parameters \(\bar p,\bar q\), has at
\(\bar q\mapsto\bar c^{\mathbf A}\) the operator, stages, ranks and stage
comparison relations of \(\varphi\).

\emph{Positive form.} Apply \cref{thm:p03-w03-p03-060-stage-comparison} to
\(\varphi'\) over \(\tau_0\): it gives LFP formulas
\(\mathrm{le}_{\varphi'}\) and \(\mathrm{lt}_{\varphi'}\) of vocabulary
\(\tau_0\) with free variables among \(\bar u,\bar w,\bar p,\bar q\). By
item~3 of \cref{thm:p03-w03-p03-059-immerman-normal-form} there are
\(\mathrm{LFP}_1\)-formulas \(\ell,\ell'\) of vocabulary \(\tau_0\)
equivalent to them on every finite \(\tau_0\)-structure. If \(\ell\) has a
free variable outside \(\bar u,\bar w,\bar p,\bar q\), prefix an
existential quantifier for it: \(\ell\) is equivalent to
\(\mathrm{le}_{\varphi'}\) under every assignment, so its truth does not
depend on the value of that variable, universes are nonempty, and
\(\mathrm{LFP}_1\) is closed under \(\exists\); likewise for \(\ell'\). Put
\[
 \mathrm{le}^+_\varphi\defeq\exists q_1\cdots\exists q_s\Bigl(\textstyle\bigwedge_{e\le s}q_e=c_e\wedge\ell\Bigr),\qquad
 \mathrm{lt}^+_\varphi\defeq\exists q_1\cdots\exists q_s\Bigl(\textstyle\bigwedge_{e\le s}q_e=c_e\wedge\ell'\Bigr)
\]
(\(\mathrm{le}^+_\varphi\defeq\ell\), \(\mathrm{lt}^+_\varphi\defeq\ell'\)
if \(s=0\)). They are formed from \(\mathrm{LFP}_1\)-formulas and
first-order atoms by \(\wedge\) and \(\exists\), so they are
\(\mathrm{LFP}_1\)-formulas of vocabulary \(\tau\), with free variables
among \(\bar u,\bar w,\bar p\). On a finite \(\tau\)-structure
\(\mathbf A\), \(\mathrm{le}^+_\varphi\) holds at \(s\) iff \(\ell\) holds in
\(\mathbf A{\restriction}\tau_0\) at \(s[\bar q\mapsto\bar c^{\mathbf A}]\),
iff \(\mathrm{le}_{\varphi'}\) holds there, iff the values of
\((\bar u,\bar w)\) are in \(\le^{\mathrm{inf}}_{\varphi'}\) for the
parameters \((s(\bar p),\bar c^{\mathbf A})\), which is
\(\le^{\mathrm{inf}}_\varphi\) for the parameters \(s(\bar p)\). The same
holds for \(\mathrm{lt}^+_\varphi\) and \(\prec^{\mathrm{inf}}_\varphi\).
\end{proof}
\end{annexitem}

\begin{annexitem}{annex:p03-w03-p03-060-gurevich-shelah}{thm:p03-w03-p03-060-gurevich-shelah}{Gurevich and Shelah: IFP collapses to LFP on finite structures}
\begin{proof}
By induction on the IFP formula \(\varphi\) we construct
\(\varphi^\dagger\) with the stated properties. Atoms are their own
translations. For \(\neg\), \(\wedge\), \(\vee\), \(\exists\), \(\forall\),
apply the constructor to the translations; the truth value is computed
from those of the immediate subformulas at the same structure and at
assignments that differ only on the quantified variable, so the induction
hypothesis gives equivalence, and no free variable is added.

Let \(\varphi=[\mathbf{ifp}\,X\bar x\,.\,\psi](\bar t)\) with \(X\) of
arity \(r\). By the induction hypothesis \(\psi^\dagger\) is an LFP formula
in which \(X\) is a free relation variable, with free first-order
variables among those of \(\psi\), and it agrees with \(\psi\) on every
finite structure under every interpretation of \(X\) and of the other free
relation variables. Hence, on a finite structure and for fixed
parameters, \(\psi\) and \(\psi^\dagger\) define the same operator, the
same inflationary stages and the same inflationary fixed point.
By the relation-binder renaming in Step~0 of the stage-comparison proof,
we may assume that \(X\) is not rebound inside \(\psi^\dagger\); only
its free occurrences are to be replaced in the two cases below.

\emph{Case \(r=0\).} \(A^0\) has the single element \(()\), and the first
stage is \(F(\varnothing)\). If \(()\in F(\varnothing)\), every stage from
\(1\) on is \(\{()\}\); otherwise \(S^1=\varnothing=S^0\) and, by (R1) in
the proof of \cref{thm:p03-w03-p03-060-stage-comparison} (which does not
use \(r\ge1\)), every stage is empty. So \(\varphi\) holds iff \(\psi\)
holds with \(X\) empty. Let \(\varphi^\dagger\) be \(\psi^\dagger\) with
every free occurrence of \(X()\) replaced by the sentence \(\bot\) of that proof; by
Replacement (with \(V=\emptyset\)) it holds iff \(\psi^\dagger\) holds with
\(X\) empty.

\emph{Case \(r\ge1\).} Choose \(r\)-tuples \(\bar u,\bar w\) of new
variables and let \(\delta_L^\ast\) be the formula built in Step~6 of the
proof of \cref{thm:p03-w03-p03-060-stage-comparison} for
\(\psi^\dagger\), after the renaming of its Step~0. Put
\[
 \varphi^\dagger\defeq[\mathbf{lfp}\,L\bar u\bar w\,.\,\delta_L^\ast](\bar t,\bar t).
\]
On a finite structure, at an assignment \(s\), it holds iff
\((\bar t^{\mathbf A}[s],\bar t^{\mathbf A}[s])\in\lambda\), that is iff
\(|\bar t^{\mathbf A}[s]|_{\psi^\dagger}<\infty\), that is iff
\(\bar t^{\mathbf A}[s]\) lies in the inflationary fixed point of
\(\psi^\dagger\), which is that of \(\psi\); this is the truth value of
\(\varphi\). The free first-order variables of \(\varphi^\dagger\) are those
of \(\bar t\) and the parameters of \(\psi^\dagger\), which are among the
free variables of \(\psi\) outside \(\bar x\); all of these are free in
\(\varphi\). Its free relation variables are those of \(\psi^\dagger\)
other than \(X\), that is those of \(\varphi\).

For the last sentence of the theorem: by
\cref{thm:p03-w03-p03-060-fp-comparison} every LFP formula is equivalent on
finite structures to an IFP formula, and by the construction just given
every IFP formula is equivalent on finite structures to an LFP formula.
\end{proof}
\end{annexitem}

\begin{annexitem}{annex:p03-w03-p03-060-av-easy}{thm:p03-w03-p03-060-av-easy}{Equivalence of the fixed-point logics forces \(\mathsf P=\mathsf{PSPACE}\)}
\begin{proof}
\emph{(E\(_{\mathrm{LFP}}\)) implies (E\(_{\mathrm{IFP}}\)).} Given
\(\psi\), take the LFP formula \(\chi\) of (E\(_{\mathrm{LFP}}\)) and apply
the translation of LFP into IFP in the proof of
\cref{thm:p03-w03-p03-060-fp-comparison}. That translation replaces
\(\mathbf{gfp}\) subformulas by negated \(\mathbf{lfp}\) formulas with the
same terms and \(\mathbf{lfp}\) by \(\mathbf{ifp}\) with the same terms, so
it adds no free variables, and the result is equivalent to \(\chi\), hence
to \(\psi\), on all finite structures.

\emph{(E\(_{\mathrm{IFP}}\)) implies (E\(_{\mathrm{LFP}}\)).} Given
\(\psi\), take the IFP formula of (E\(_{\mathrm{IFP}}\)) and apply
\cref{thm:p03-w03-p03-060-gurevich-shelah}, which adds no free variables.

\emph{(E\(_{\mathrm{LFP}}\)) implies \(\mathsf P=\mathsf{PSPACE}\).} We
verify item~2 of \cref{thm:p03-w03-p03-059-ptime-pspace}. Let \(\tau\) be a
finite relational vocabulary and \(\mathcal K\) a class of ordered finite
\(\tau\)-structures. If \(\mathcal K\) is defined among them by a PFP
sentence \(\psi\) of vocabulary \(\tau\cup\{<\}\), then
(E\(_{\mathrm{LFP}}\)) with \(\sigma=\tau\cup\{<\}\) gives an LFP formula
\(\chi\) without free variables, that is a sentence, equivalent to
\(\psi\) on all finite \((\tau\cup\{<\})\)-structures, in particular on the
ordered ones; so \(\chi\) defines \(\mathcal K\). If \(\mathcal K\) is
defined by an LFP sentence, then by \cref{thm:p03-w03-p03-059-lfp-in-pfp}
it is defined by a PFP sentence. So item~2 of
\cref{thm:p03-w03-p03-059-ptime-pspace} holds, and that corollary gives
\(\mathsf P=\mathsf{PSPACE}\). Through the equivalence just proved,
(E\(_{\mathrm{IFP}}\)) implies it as well.
\end{proof}
\end{annexitem}
